\documentclass[11pt,a4paper]{article}

% -------------------------------------------------------
% Page layout and typography
% -------------------------------------------------------
\usepackage[margin=2.5cm]{geometry}
\usepackage[T1]{fontenc}
\usepackage[utf8]{inputenc}
\usepackage{lmodern}
\usepackage{microtype}
\usepackage{setspace}
\onehalfspacing

% -------------------------------------------------------
% Tables, figures and maths
% -------------------------------------------------------
\usepackage{graphicx}
\usepackage{float}
\usepackage{booktabs}
\usepackage{tabularx}
\usepackage{longtable}
\usepackage{array}
\usepackage{amsmath}
\usepackage{siunitx}
\usepackage{gensymb}

% -------------------------------------------------------
% TikZ figures
% -------------------------------------------------------
\usepackage{tikz}
\usetikzlibrary{
    arrows.meta,
    positioning,
    shapes.geometric,
    shapes.misc,
    calc,
    fit,
    backgrounds
}

% -------------------------------------------------------
% References and links
% -------------------------------------------------------
\usepackage[round,authoryear]{natbib}
\usepackage[hidelinks]{hyperref}
\usepackage[nameinlink,noabbrev]{cleveref}

% -------------------------------------------------------
% Metadata
% -------------------------------------------------------
\title{
What Makes Mountain Wildfire Different?\\
Evidence from the European Alps and Other Ranges
}

\author{
% DRAFT author list: order, institutes and e-mail to be confirmed by the author team
Alexandre Dunant$^{1,*}$, Alice Crespi$^{1}$, Michael de Rachewiltz$^{1}$, Christian Hoffmann$^{1}$,\\
Marco Mina$^{1}$, Massimiliano Pittore$^{1}$, Uta Schirpke$^{1}$, Stefano Terzi$^{1}$,\\
Jenny Ufer$^{1}$, Marc Zebisch$^{1}$ \\
\small $^{1}$Eurac Research, Bolzano, Italy \\
\small $^{*}$Corresponding author: [e-mail to be added]
}

\date{}

% -------------------------------------------------------
% Begin document
% -------------------------------------------------------
\begin{document}

\maketitle


\begin{abstract}
\noindent
Wildfire research in Europe has concentrated on Mediterranean regions, but changing fire weather and land use are increasing concern in the European Alps. We review Alpine evidence alongside studies from other mountain ranges to assess how terrain, forest protection and management affect wildfire risk. Most recorded Alpine fires are small, yet a fire above a settlement or transport corridor can have consequences beyond its perimeter. Loss of vegetation and changes in soil condition can increase erosion, debris-flow susceptibility and rockfall exposure. These effects recover at different rates: soil changes can persist for years, while reduced forest protection may last for decades. Mountain weather and heterogeneous fuels also complicate the transfer of danger indices and spread models. The evidence supports local evaluation of these tools, combined with improved records of small fires and post-fire processes. Forest and land-use management are central to prevention, but treatment must account for protective function, access and species response. Comparative studies identify relevant mechanisms and management options without establishing that their numerical effects transfer to the Alps. We conclude that mountain wildfire assessment needs to connect fire likelihood and behaviour with exposed assets, vulnerability and the recovery of forest protection. Priorities include sustained post-fire monitoring, evaluation of treatment trade-offs and coordination between fire, forest and natural-hazard authorities.
\end{abstract}

\noindent\textbf{Keywords:}
mountain wildfire; European Alps; mountain and lowland fire regimes; fire regime change; protection forest; fuel management; drought and water availability; post-fire hazards; wildfire governance and economics; risk assessment


% =======================================================
% =======================================================
\section{Introduction}
\label{sec:intro}

\subsection{An emerging fire frontier in temperate mountains}

The geography of European fire science has followed the geography of burnt area. Because the Mediterranean basin accounts for most of the continent's burning, it has also supplied most of the continent's fire statistics, operational danger rating and prevention doctrine. That arrangement is now being unsettled by a change in where fire occurs. Burnt area has been declining in the Mediterranean while increasing in mid- and high-latitude regions where fuel is abundant and where plant communities carry little adaptation to fire \citep{miller_compounding_2025}, and the shift is vertical as well as latitudinal. Across eight European mountain regions between 2000 and 2025, the median elevation of fires larger than 30\,ha rose from 910 to 1004\,m and the ninetieth percentile from 1353 to 1460\,m, an upslope migration of approximately 72\,m per decade driven principally by vapour pressure deficit \citep{beloiu_climatedriven_2026}.

This pattern is not confined to mountains. In the Netherlands, a densely populated temperate country with no fire tradition, 3{,}667 landscape fires were recorded between 2017 and 2022, averaging 611 a year and around 405\,ha annually. Mean fire size was 1.5\,ha, 88\% of events burned less than a hectare, and almost all were human-ignited \citep{stoof_temperate_2024}. That profile is close to the Alpine one. So is the institutional position it describes, in which fire services are specialised in urban firefighting and neighbouring countries hold no systematic wildfire statistics at all.

Projections for the Alps themselves are large relative to that baseline. Hazard probability across the range is projected to increase by 15.5\% under SSP1-2.6 and by 35.4\% under SSP5-8.5, with the largest relative increases in the northern and eastern sectors rather than in the already fire-affected south-west \citep{gerberding_mapping_2025}. In adjacent central European landscapes, a regional large ensemble under RCP8.5 places the emergence of the forced fire-danger signal before 2050 in every subregion examined, and shifts present-day hundred-year fire weather to a ten-year recurrence by 2090 \citep{miller_climate_2024}. Both results describe a hazard moving from the anecdotal to the routine within the planning horizon of the institutions that would have to absorb it.

How much of that increase is realised is not, however, fixed by climate alone. Coupling two fire models, SPITFIRE and the Burnt Area Simulator for Europe, to the fire-enabled dynamic global vegetation model LPJmL, intensified fire weather alone is projected to raise annual burned area by around 39\% under SSP1-2.6 and by nearly 192\% under SSP3-7.0 by the end of the century. Continued improvement in management capacity would reduce burned area by 72\% to 92\% relative to scenarios without it, while still leaving net increases in roughly 55\% of Europe's fire-prone regions, with the hotspots under the higher pathway concentrated in foothills and medium-elevation terrain including the Alps, Pyrenees, Carpathians and Massif Central \citep{billing_future_2026}. Management capacity is therefore a first-order determinant of Alpine outcomes rather than a residual after the climate signal has been counted.

Meanwhile, the wider European context has moved. Burnt area and fire carbon emissions across Europe reached record levels in 2025, concentrated in the Iberian Peninsula, and the same assessment finds that emissions mitigation has already reduced projected end-of-century burnt area in extreme events, by 4\% to 23\% under a high pathway relative to a worst case now considered implausible, and by a further 5\% to 36\% under a Paris-consistent pathway \citep{burton_state_2026}. Alpine fire danger is thus conditional on mitigation as well as on adaptation, which is not how it is usually presented.

Fire is therefore arriving in terrain where neither the science, the management apparatus, nor the exposed population has accumulated experience of it. The deficit was named more than a decade ago from inside Alpine forestry. Reviewing the risks facing Alpine Space forests, the MANFRED consortium observed that a history of active suppression meant fire risk did not receive full attention in the Alps, and suggested that future climate conditions might well require an improved alerting system \citep{jandl_management_2013}. The intervening decade has supplied the climate signal. The attention has followed more slowly.

\subsection{Why mountain fire is a distinct problem}

The Mediterranean toolkit does not transfer cleanly to mountains, and the reasons begin with the atmosphere. The assumption of a continuous fire-weather field, interpolable from valley-floor stations, loses validity where upper winds interact with complex relief to produce discontinuous and locally amplified surface conditions \citep{sharples_overview_2009}. The magnitude of the effect has been quantified. Atmospheric profiles capable of generating enhanced surface winds occurred in 6.5\% of cases over the least rugged terrain of a Tasmanian study domain but in 40.2\% over the most rugged, and fire simulations driven by surface winds alone degraded systematically as terrain complexity increased \citep{ozaki_impact_2022}.

The same relief that distorts the wind field also reorganises the behaviour of the fire it drives. Rate of spread accelerates non-linearly with upslope angle, roughly two and a half times between 20\degree{} and 30\degree{} in controlled experiments, and the dominant preheating mechanism switches from radiation to convection across that range \citep{hayajneh_wind_2026}. The landscape-scale consequence is measurable. Across 21 Swiss pyroregions, the length of upslope paths within forested area, combined with the aggregation of that area, explains between 56\% and 65\% of the variance in large-fire statistics without any weather term \citep{conedera_linking_2024}.

Weather and spread would alone justify a separate treatment, but the deeper difference lies in what happens after the fire is out, because in mountains the fire is rarely the end of the event. Alpine forests are managed substantially for their protective function against rockfall, avalanche and shallow landslide, and fire removes that function over a window measured in years \citep{melzner_wildfire-induced_2025, melzner_brief_2019}. The asset at risk is typically not the stand that burns but the settlement, road or railway beneath it. Each of these departures recurs in a different part of what follows, and together they explain why a burnt-area statistic, which serves the Mediterranean reasonably well, says little about mountain consequence.

\subsection{Objective, argument and the structure of this review}

The Alpine fire community already has a consensus statement. The EUSALP white paper of February 2020 set out the state of knowledge for policy makers, established the arc-wide statistical baseline, identified the loss of protective function as the central concern, and proposed a costed integrated fire management framework \citep{muller_forest_2020}. We take that position as our point of departure and ask what the evidence has since established. In several places that evidence complicates the 2020 position rather than confirming it, and where it does we identify the finding responsible.

Here we assemble the evidence needed to treat mountain wildfire as a systemic risk rather than as a forest-loss problem. Part of that argument extends the 2020 position and part of it complicates it, and we keep the two apart because an explicit account of what is inherited is more useful than implied originality.

The extension begins from framings already established in 2020. The Alpine fire regime is structurally distinct rather than merely smaller than the regimes on which European fire science was built. That claim was advanced in 2020 from the arc-wide statistics, and subsequent work has made it quantitative in the vertical dimension, through a measured upslope migration of activity in European mountains \citep{beloiu_climatedriven_2026}, and in the operational one, through terrain-conditioned estimates of the difficulty of containment \citep{cardil_assessing_2025}. Similarly, because what fire destroys in the Alps is a managed protective function, the response is substantially a forestry problem. That identification is also inherited. The material for acting on it is not, and we therefore treat silviculture, fuel treatment and post-fire restoration as first-order subjects on the strength of monetised valuations of the protective service \citep{accastello_asforesee_2019, moos_climate_2021}, a controlled Alpine prescribed-burning trial \citep{gamba_prescribed_2025} and severity-resolved measurement of the interval over which protection is lost \citep{melzner_wildfire-induced_2025}.

The complications are the substance of the review, and the most far-reaching concerns consequence. Consequence in the Alps is decoupled from burnt area and runs through several pathways at once, namely direct loss of protective structure, ecological and hydrological change, and a geomorphic cascade in which fire sets a susceptibility that a later hydro-meteorological event realises. The geomorphic pathway, although the least represented in current assessment, should be understood as one term in the loss function rather than as the whole of it. Part of that pathway also falls outside the return-period logic in which mountain hazard is conventionally expressed, since a burnt Alpine catchment has been observed to respond to a rainfall intensity below the threshold that unburnt basins of its type require \citep{laporta_postwildfire_2025}.

If consequence escapes the burnt-area metric, it also escapes the apparatus built on that metric, which carries hazard some way towards risk but stops short of it. The exposure and vulnerability terms that would convert a burn probability into an expected loss exist as instruments but have not been assembled for the arc. The Alpine data record is meanwhile too thin to justify the conventional remedy of added complexity, and we use that constraint, rather than any appeal to greater model sophistication, to organise a research agenda.

The status of the 2020 statement gives the exercise a further purpose. It is a policy white paper rather than a contribution to the scientific record. Its evidence base is not indexed, peer reviewed or available for citation and appraisal in the ordinary way, and its positions cannot readily be tested or extended by the community that would have to act on them. Setting that position against the peer-reviewed literature that has since accumulated, and stating for each element how securely it is now supported, is itself part of what a review can supply.

We organise the material by the domains in which the evidence is produced rather than by the stages of a fire. Section~\ref{sec:hazard} treats wildfire as a natural hazard, covering initiation, spread, hazard modelling, and the datasets and Earth observation on which all of it depends. Section~\ref{sec:climate} treats the climatic and meteorological drivers, including the wind systems, lightning climatology and projections specific to complex terrain. Section~\ref{sec:water} follows the water pathway, from drought and fuel moisture through soil alteration and the geomorphic cascade to water supply, snow and ice. Section~\ref{sec:forest} treats the forest itself, its recent change, its ecology and its management. Section~\ref{sec:risk} assembles the exposure and vulnerability terms that convert hazard into risk, and Section~\ref{sec:governance} treats governance, policy and economics. Each of these sections closes with an explicit comparison between the mountain case and the non-mountain baseline, and we collect those comparisons in Table~\ref{tab:contrast}. We close by considering what the observational record can support.

% =======================================================
\section{Review Design and Scope}
\label{sec:scope}

\subsection{Scope and inclusion rule}

Our question is what distinguishes wildfire in mountains from wildfire in the lowland and Mediterranean settings on which European fire science was built. The European Alps are the focal case, because they carry the densest and most recent evidence and because the practitioner and policy community is organised at that scale, but the claim being tested is about mountains. The corpus is therefore built in two tiers rather than one core and a set of exceptions.

Tier one is the Alpine arc as delimited by the Alpine Convention, covering Austria, France, Germany, Italy, Liechtenstein, Monaco, Slovenia and Switzerland. This is the perimeter within which harmonised Alpine fire statistics were first assembled \citep{valese_wildfires_2011, muller_forest_2020}, within which arc-scale hazard mapping has been performed \citep{gerberding_mapping_2025}, and across which the practitioner community is organised \citep{hoang_stakeholders_2026}.

Tier two is mountain fire systems generally, European and otherwise, admitted on the strength of the process rather than by exception. A result from the Pyrenees, the Carpathians, the Apennines, the Scandes, the Rocky Mountains or the Sierra Nevada enters wherever it bears on a process the review treats, whether or not Alpine evidence on that process exists. Where both are available the comparison is the point: agreement across ranges indicates that a departure belongs to mountains rather than to the Alps, and disagreement locates what is specific to the Alpine setting. Table~\ref{tab:ranges} characterises the systems drawn on and what each contributes.

\begin{table}[H]
\centering
\footnotesize
\setlength{\tabcolsep}{4pt}
\renewcommand{\arraystretch}{1.1}
\caption{Mountain systems drawn on in this review, with the regime feature that makes each informative and the contribution it makes to the mountain-versus-lowland comparison.}
\label{tab:ranges}
\begin{tabularx}{\textwidth}{@{}>{\raggedright\arraybackslash}p{2.1cm} >{\raggedright\arraybackslash}p{3.0cm} >{\raggedright\arraybackslash}X >{\raggedright\arraybackslash}p{2.9cm}@{}}
\toprule
\textbf{Range or study region} & \textbf{Evidence used} & \textbf{Contribution to the comparison} & \textbf{References} \\
\midrule
European Alps & Seasonal ignition regimes and forests protecting downslope assets & Focal evidence on fire occurrence, protective-function loss and post-fire hazards & \citep{conedera_characterizing_2018, muller_forest_2020, gehring_shallow_2019} \\
Pyrenees (Aran valley) & Winter burning and participatory fire planning in an Atlantic mountain valley & One operational precedent for managed fire; not a description of the whole range & \citep{oliveres_path_2025} \\
Carpathians & Spring ignition peak and greater susceptibility at lower elevations & Links between land use, access and the spatial distribution of fire & \citep{manczinger_fire_2025} \\
Northern Apennines & Long-term replacement of mixed fir forest by beech & Land-use and fire legacies in present forest composition & \citep{moralesmolino_longterm_2021} \\
Dinarides; Tatras & Lightning-related infrastructure exposure; post-windthrow salvage & Separate comparisons for infrastructure and disturbance legacies & \citep{kyriakopoulos_impact_2023, seben_impact_2024} \\
Fennoscandia & National records detect many fires absent from satellite products & Detection limits in a northern region containing both mountain and lowland settings & \citep{bakke_datadriven_2023} \\
Rocky Mountains & Post-fire debris-flow monitoring and terrain-related fire refugia & Rainfall thresholds, model validation and topographic controls & \citep{rengers_evaluating_2024, rogeau_effect_2018} \\
Sierra Nevada; western US & Observed changes in snow cover and melt after fire & Cryospheric processes for which Alpine measurements remain limited & \citep{micheletty_application_2014, koshkin_wildfire_2022} \\
Southern Appalachians & Wildfire after fire exclusion; soil microbial response & Disturbance history and ecological endpoints for treatment evaluation & \citep{reilly_drivers_2022, snyder_taxonomic_2026} \\
Australian Alps; Tasmania & Mountain governance and terrain-modified winds & Separate institutional and atmospheric comparisons & \citep{lockwood_biodiversity_2014, ozaki_impact_2022} \\
East Asian study regions & Mountain-valley winds and slope-dependent smouldering experiments & Mechanistic evidence; experimental effects are not Alpine parameter estimates & \citep{wang_dynamic_2026, shan_influence_2024} \\
\bottomrule
\end{tabularx}
\end{table}

Drawing on other ranges requires discipline about what travels. Mechanisms are treated as transferable where the physics is general, coefficients are not, and the distinction is marked at the point of use. A result from outside the Alps accordingly serves to establish that a process operates, to bound its magnitude, or to show that an Alpine observation is or is not exceptional, while Alpine parameters are estimated from Alpine data alone.

The comparison addresses processes shared across mountain environments: terrain-modified winds, slope-dependent spread, sediment mobilisation and loss of forest protection. These processes also occur outside the Alps, although their importance varies among settings. Contextual literature on mountain processes without fire, and on fire outside mountains, supplies definitions, methods and comparison cases. It is distinguished from direct evidence of Alpine fire effects.

\subsection{Literature selection, review and evidence weighting}

We combined keyword searches, citation searches and thematic appraisal to assemble the literature for this review. Figure~\ref{fig:flow} shows the sequence from the starting literature to disciplinary review. Appendix~\ref{app:protocol} distinguishes the publications used to initiate the search, the candidate records screened and the sources supporting the synthesis.

The starting literature comprised the EUSALP white paper \emph{Forest fires in the Alps} \citep{muller_forest_2020} and \emph{Mapping the probability of forest fire hazard across the European Alps under climate change scenarios} \citep{gerberding_mapping_2025}. These provided the regional management context and a spatial assessment of future hazard, respectively (Step~1).

We searched Google Scholar using the combinations \emph{wildfire}, \emph{mountain} and \emph{Alps}, then \emph{wildfire} and \emph{mountains} to include other ranges (Step~2). Selection followed the thematic scope of Sections~\ref{sec:hazard}--\ref{sec:governance}, informed by the disciplines represented in the author team. The initial library contained 49 distinct publications. This is the starting set for citation searches, rather than the size of the final evidence base.

Citation searches identified additional candidate publications from the reference lists of the starting literature and from works citing 14 core Alpine publications (Step~3). Bibliographic records were matched by DOI or normalised title to remove repeated entries. The output of this stage was a pool of potentially relevant literature; identification alone did not establish eligibility for the synthesis.

Targeted searches and author-supplied material supplemented topics that remained insufficiently covered during drafting (Step~4). These additions were assessed against the same scope and evidence requirements as the initial candidates. Availability of full text constrained this stage, a limitation discussed in Section~\ref{sec:discussion}.

Screening classified records as Alpine evidence, comparative mountain evidence, context or outside scope, with a reason recorded for each decision (Step~5). The register contains 3{,}228 screened records, of which 896 were retained as relevant candidates and 2{,}332 excluded. Retention in this register does not imply that every candidate contributes to the paper. For sources used in the synthesis, evidence notes link the study setting, methods, results and limitations to the claims supported. This distinction separates the search record from the evidence presented in the review.

Disciplinary review by the co-authors checks coverage, interpretation and the strength of the conclusions in each thematic section (Step~6). This stage remains open in the present draft. Papers suggested during review return to Step~3, where their references and citing literature can extend the search. The resulting candidates then pass through the same screening and appraisal stages before incorporation. The return arrow in Figure~\ref{fig:flow} denotes this procedure; it does not represent a completed additional search round.

\begin{figure}[H]
\centering
\begin{tikzpicture}[
  scale=0.8, transform shape,
  font=\small,
  stp/.style={draw, rounded corners=2pt, align=left, text width=6.4cm, inner sep=5pt, fill=black!6},
  res/.style={draw, rounded corners=2pt, align=left, text width=6.2cm, inner sep=5pt},
  arr/.style={-{Latex[length=2mm]}, thick},
  lnk/.style={-{Latex[length=2mm]}, thin}
]
\node[stp] (s1) {\textbf{Step 1. Starting literature}\\ Alpine fire management white paper \citep{muller_forest_2020}; Alpine hazard mapping under climate scenarios \citep{gerberding_mapping_2025}};
\node[stp, below=5mm of s1] (s2) {\textbf{Step 2. Keyword searches}\\ Google Scholar: \emph{wildfire} + \emph{mountain} + \emph{Alps}; then \emph{wildfire} + \emph{mountains}};
\node[stp, below=5mm of s2] (s3) {\textbf{Step 3. Citation searches}\\ Reference lists of the starting literature and publications citing core Alpine studies};
\node[stp, below=5mm of s3] (s4) {\textbf{Step 4. Thematic supplementation}\\ Targeted searches and author-supplied papers addressing gaps in coverage};
\node[stp, below=5mm of s4] (s5) {\textbf{Step 5. Screening and appraisal}\\ Relevance assessment, source interpretation and documentation of supporting evidence};
\node[stp, below=5mm of s5, dashed] (s6) {\textbf{Step 6. Disciplinary review}\\ Co-author assessment of coverage, interpretation and conclusions};

\node[res, right=8mm of s1] (o1) {Regional context for the review: observed fire regimes, management priorities and projected hazard};
\node[res, right=8mm of s2] (o2) {Initial library: 49 distinct publications used to begin citation searches};
\node[res, right=8mm of s3] (o3) {Candidate literature for screening: 3{,}375 distinct works from reference lists; 319 citing works with identifiable bibliographic details. Overlap removed before screening};
\node[res, right=8mm of s4] (o4) {Additional candidate papers on hazard, climate, water, forests, risk and governance};
\node[res, right=8mm of s5] (o5) {Screening register: 3{,}228 records assessed, 896 relevant candidates retained, 2{,}332 excluded. Evidence notes supporting the thematic synthesis};
\node[res, right=8mm of s6, dashed] (o6) {Revised interpretation and coverage; additional papers proposed for assessment. Review remains open};

\foreach \i/\j in {s1/s2, s2/s3, s3/s4, s4/s5, s5/s6} \draw[arr] (\i) -- (\j);
\foreach \i/\j in {s1/o1, s2/o2, s3/o3, s4/o4, s5/o5, s6/o6} \draw[lnk] (\i) -- (\j);
\draw[arr] (s6.west) -- ++(-8mm,0) |- (s3.west);
\node[font=\scriptsize, rotate=90, anchor=south] at ([xshift=-9mm]s4.west) {Additional papers from review};
\node[font=\footnotesize\bfseries, anchor=south] at ([yshift=1mm]s1.north) {Review stage};
\node[font=\footnotesize\bfseries, anchor=south] at ([yshift=1mm]o1.center |- s1.north) {Contribution to the review};
\end{tikzpicture}
\caption{Literature selection and review. The main arrows show the sequence from starting literature to disciplinary review. The return arrow allows papers suggested during review to initiate further citation searches and pass through the same screening process. Dashed outlines identify the review stage that remains open. Counts describe the current search register, not the number of studies cited; Appendix~\ref{app:protocol} gives the detailed accounting.}
\label{fig:flow}
\end{figure}

Our synthesis is thematic rather than meta-analytic, and this follows from the evidence rather than from preference. Alpine fire studies differ in their definitions of severity, in burnt-area thresholds, in reporting periods and in the administrative basis of their fire records, so effect sizes are not commensurable across studies. The same obstacle has been documented for the neighbouring literature on land abandonment and disturbance, where study sites cluster strongly and coverage across disturbance types is uneven \citep{mantero_influence_2020}.

Much of the Alpine evidence base consists of individual documented events, and we therefore adopt an explicit convention for weighting it. We use case evidence to establish that a process operates, to characterise its mechanism and to bound its magnitude in one setting, reserving regional parameters for multi-site evidence, and we state the number of independent cases wherever a claim rests on fewer than three. Where a mechanism is established outside Europe, we identify it as imported and use it for mechanism only. Where a claim is supported by multiple independent Alpine studies, we say so. Single-case magnitudes should accordingly be read as existence proofs with an observed value attached rather than as central estimates, and several of the quantities cited in Sections~\ref{sec:water} and~\ref{sec:forest} fall into this category.

% =======================================================
\section{Wildfire as a Natural Hazard}
\label{sec:hazard}

\subsection{Initiation: the Alpine ignition regime and its modelling}

The regime that any Alpine ignition model must reproduce is defined less by its magnitude than by its shape. Harmonised statistics for 2006--2017 place annual fire numbers between 37 in Slovenia and 1043 in Italy, and annual burnt area between 64\,ha in Austria and 9984\,ha in Italy; normalised by forest area, the range spans two orders of magnitude, from 0.22 to 44.13\,ha burnt per 10{,}000\,ha of forest per year \citep{muller_forest_2020}. Roughly 80\% of Alpine fires burn less than one hectare, and median fire size falls below 0.1\,ha in Austria, Germany and Switzerland. The distribution is nonetheless strongly heavy-tailed: in the French and Italian Alps the mean fire covers about a hectare, yet the 10\% of fires exceeding ten hectares account for some 85\% of total burnt area \citep{muller_forest_2020}. Both halves of that statistic matter, and citing either alone misleads. The many small fires define what monitoring must detect and what suppression routinely handles; the few large ones define the burnt area, and, in an Alpine setting, the events capable of removing protective function across a whole slope. For comparison, the mean patch size of all detected forest disturbance across Europe over three decades is 1.09\,ha \citep{senf_mapping_2021}, so Alpine fire sits at the fine end of an already fine-grained disturbance regime.

Seasonality is bimodal, and we take this to be the regime's most consequential structural feature. A spring season of fast-spreading surface fires at low elevation peaks in March and April; a summer season peaks in July and August; and winter fires occur in foehn-influenced valleys \citep{valese_wildfires_2011, muller_forest_2020}. The pattern is not confined to the Alpine arc: in the Czech Republic the same bimodality appears with an April peak and a July--August peak, the former attributed to forestry operations resuming after winter and the latter to heat, drought and tourism, with a threshold response above 18\,\degree C \citep{bercak_combination_2024}. The two modes have different controls, and there is no reason to expect them to change in step.

Human causes account for approximately 90\% of recorded Alpine fires, while lightning contributes around 10\% annually and up to half of summer ignitions in some areas \citep{muller_forest_2020, valese_wildfires_2011}. In Austria, access networks and agricultural and forestry activity help explain human ignition patterns \citep{arndt_modeling_2013}. However, access is not a simple proxy for ignition pressure. In the Valais, fire occurrence increased with road density up to about 30\,m per hectare and then declined, possibly reflecting intensive land use, reduced fuels and easier suppression \citep{zumbrunnen_human_2012}. Roads and settlements also covary with elevation and temperature in steep terrain, complicating attribution to human activity alone \citep{bekar_crossregional_2020}.

Ignition controls also differ by season. In one inner-Alpine valley, summer ignition was associated with maximum temperature and precipitation seasonality, while winter ignition was associated with precipitation and elevation. Proximity to infrastructure mattered in winter but not in summer \citep{vacchiano_modeling_2018}. These differences support seasonal modelling rather than an annual relationship that averages across distinct processes.

A spring peak also occurs in the Carpathians. A susceptibility analysis covering seven countries and 5{,}173 fires during 2010--2020 identified cropland and pasture among the leading predictors, with a March maximum and greater susceptibility at lower elevations \citep{manczinger_fire_2025}. The reported association with easier access on flatter ground illustrates how terrain and human activity interact. It does not contradict the upward trend over time reported for large fires in European mountains: the studies measure different aspects of fire occurrence.

Year-to-year departures from that seasonal structure are large enough to be diagnostic rather than incidental. Between 1993 and 2013 Austria recorded three consecutive anomalous seasons in different parts of the calendar: autumn 2011 with 47 fires and 27\,ha against a mean of 11 and 3; spring 2012 with 184 fires, the highest spring count in the series, against a mean of 57; and summer 2013 with 107\,ha against a mean of 22. Comparison across the arc places the larger anomalies in the eastern Alpine countries, Austria, Slovenia and northeastern Italy, rather than in the west \citep{muller_anomalies_2015}. A regime whose interannual variance is this high is one in which a short record cannot separate trend from fluctuation, and the same study reports that roughly half of Austrian fire records lack a burnt area at all.

Beneath this interannual noise the regime regionalises coherently, with Southern, Northern and Maritime clusters separating the high fire density of the southern slope from the substantially lower burnt-area fraction of the north and differing in drought frequency and duration, vegetation type, mean elevation, valley orientation, population density and educational level \citep{conedera_characterizing_2018}. With only ten years of record the classification functions as a baseline against which future regime shifts and anomalous seasons can be detected rather than as a description of a stable state, and on present evidence fire does not rival storms and bark beetles for damage in the Alpine region, although future regimes may evolve under climatic and socio-economic change \citep{conedera_characterizing_2018}. That is a more measured framing than much of the recent literature adopts, and this review adopts it.

Ignition records support several statistical approaches. In Trentino--South Tyrol, a binomial generalised additive model (GAM) combines static landscape predictors with time-varying precipitation and temperature to estimate ignition probability and operational thresholds \citep{moreno_1_nodate}. In a separate European study, pooling regions reduced local discrimination relative to site-specific fitting \citep{gelabert_assessing_2025}. However, Maxent models fitted across southern Switzerland, Carinthia, Valencia and south-western Turkey were more robust in transfer than models trained in individual regions \citep{bekar_crossregional_2020}. Local discrimination and transfer performance are therefore distinct criteria. Moving from 1\,km to 100\,m resolution also improved performance in the latter study by better representing local landscape structure. Neither test establishes performance under future climatic conditions.

\subsection{Spread: fire behaviour and spread modelling}

On spread, the Alpine observational record and the process literature are very unevenly developed. The observational evidence for terrain-driven behavioural anomaly in the Alps rests, at present, on a single documented case. During the 2018 Hallstatt fire in Austria, vertical rock walls, anabatic winds and patchy vegetation produced upward jumping of the fire and a spotty burn pattern, generating spatially variable intensity and correspondingly heterogeneous changes in soil and rock mass \citep{melzner_brief_2019}. That heterogeneity is not a curiosity: it propagates into the severity field and from there into post-fire hazard, which is why terrain-driven complexity in fire behaviour is a risk-assessment problem and not only an operational one.

Experimental evidence outside the Alps shows that wind and slope can change the dominant heat-transfer mechanism and produce nonlinear spread responses \citep{hayajneh_wind_2026}. Thresholds reported for controlled fuel beds depend on fuel arrangement, geometry and experimental conditions. They help explain possible behaviour in steep terrain, but should not be treated as operational thresholds for Alpine valleys.

What the Alpine literature documents case by case, coupled fire--atmosphere modelling can now resolve explicitly, and the demonstration exists for a European event outside the Alps. Simulation of the 2017 Pedr\'og\~ao Grande fire at 80\,m resolution reproduced updraughts above 15\,m\,s$^{-1}$, a convective column reaching roughly 10\,km, a plume top rising from below 4\,km to over 9\,km, and a localised downward outflow of similar magnitude developing ahead of the fire front, with wind-direction changes triggering acceleration and merger of fire fronts \citep{couto_numerical_2024}. Alpine valleys generate the terrain-forced circulations that such models are built to resolve, and no equivalent coupled simulation has been published for an Alpine fire. This is the methodological gap that separates the descriptive Alpine meteorological literature from the mechanistic one.

A comparison of FlamMap 6.2 and PROPAGATOR examined three Lombardy fires: Tirano, Campolaro and Sonico. FlamMap combines Rothermel's surface-spread formulation with a minimum-travel-time algorithm under fixed environmental conditions; PROPAGATOR uses a stochastic cellular automaton that can incorporate changing weather. Both reasonably reproduced propagation, but their intensity and spread-rate estimates differed substantially \citep{malanchini_modeling_2025}. The study favoured FlamMap for detailed planning and PROPAGATOR for flexible operational use. These results support evaluating model outputs against their intended use, rather than inferring general superiority from perimeter agreement. Fireline intensity is relevant to damage, but severity also depends on heating duration, fuels and the vulnerability of vegetation and soil.

The wider experimental literature helps explain why spread formulations can diverge in steep terrain. Interactions among wind, slope, fuel continuity and convective heating become important outside the conditions used to calibrate empirical models \citep{hayajneh_wind_2026}. Alpine evaluation therefore needs observations of spread and energy release, together with the weather and fuel inputs that determine them.

Terrain also determines where fire does not go. That side of the relationship has been quantified outside Europe. Across 912{,}000\,ha of the Canadian Rocky Mountains, 557 mapped patches of persistent fire refugia covering 27.7\% of the forested area were predicted from terrain and fuel with elevation the leading variable at 23.7\% relative influence, followed by the proportion of non-fuel in the surrounding landscape at 20.3\%, solar radiation as an aspect proxy at 14.6\% and topographic position at 10.1\% \citep{rogeau_effect_2018}. The skill is strongly scale-dependent: the area under the curve rises from 0.71 for refugia of all sizes to 0.84 for patches under 10\,ha, so terrain predicts small unburnt patches far better than large ones, and the authors note that topography may matter less during large fire events. For a range whose fire-size distribution is dominated by small events, that scale dependence is favourable, and it is the same geometric argument that supports terrain-based prioritisation in the Alps.

Coupled models can resolve what the empirical ones prescribe, and the terrain circulation they resolve is itself a spread driver. Simulation of a 2023 fire in the Greater Khingan Mountains at high resolution, reproducing the observed burnt area with a spatial similarity of 0.82, attributes its fastest phase to nocturnal down-valley winds of 3 to 5\,m\,s$^{-1}$: maximum spread rate rose from 0.3 to 1.1\,m\,s$^{-1}$ and peak heat flux from 11 to 159\,kW\,m$^{-2}$, then fell back as daytime up-valley winds reasserted themselves \citep{wang_dynamic_2026}. The setting is Chinese boreal rather than Alpine, and the coefficients are not transferable, but the mechanism is the diurnal valley circulation that operates in every Alpine valley, and it implies that the hours of fastest Alpine spread may be the hours when crews are least able to work.

Slope also governs a mode of spread that perimeter-based evaluation does not see. Alpine lightning fires are reported as harder to extinguish because they burn below the surface, and controlled smouldering experiments show that slope acts on that mode as strongly as on flaming spread: between level ground and a twenty-degree slope, peak temperature rose from 769 to 948\degree{}C and both vertical and horizontal spread rates roughly doubled, with the slope effect on temperature exceeding its effect on rate \citep{shan_influence_2024}. The experiments are Chinese boreal rather than Alpine, so the coefficients do not transfer, but the direction does, and it bears on the mop-up phase that occupies most of the effort in an Alpine incident.

Severity itself appears only partly predictable from terrain. Across 113 large Italian fires, pre-fire land cover was the strongest predictor of burn severity, with conifer stands highest, followed by elevation, while slope contributed weakly; distance to roads and railways entered positively, indicating higher uncontrolled severity in remote terrain \citep{malandra_burn_2022}. Fuel type leads, access modulates, and slope, so dominant for spread, is secondary for severity.

\subsection{Fire danger and hazard modelling}

The Canadian Fire Weather Index (FWI) is widely used in Europe, usually with locally defined danger classes. Its weather-based formulation does not explicitly represent terrain or spatial variation in fuel type \citep{zacharakis_environmental_2023}. Italy's RISICO system adds fuel information and produces hourly dead fuel moisture, spread-rate and intensity estimates on a 1\,km grid \citep{perello_fuelaware_2025}. Its fuel classification incorporates susceptibility information that includes topography. Live fuel moisture remains fixed, however, and the reported index evaluation retained fires larger than 70\,ha because of its grid and pixel-selection criteria. This limits validation for small Alpine fires; it is not a physical minimum fire size.

Because the instruments are rarely named together, Table~\ref{tab:models} sets out those that appear in the Alpine and European evidence reviewed here, with the constraint each meets in steep terrain.

\begingroup
\footnotesize
\setlength{\tabcolsep}{4pt}
\renewcommand{\arraystretch}{1.08}
\begin{longtable}{@{}>{\raggedright\arraybackslash}p{2.5cm} >{\raggedright\arraybackslash}p{3.6cm} >{\raggedright\arraybackslash}p{5.0cm} >{\raggedright\arraybackslash}p{3.5cm}@{}}
\caption{Models and indices used in the reviewed studies. Outputs and limitations refer to the cited application, rather than to every use of a model.}\label{tab:models}\\
\toprule
\textbf{Model or index} & \textbf{Method and output} & \textbf{Constraint on interpretation} & \textbf{References} \\
\midrule
\endfirsthead
\multicolumn{4}{l}{\tablename\ \thetable\ continued}\\
\toprule
\textbf{Model or index} & \textbf{Method and output} & \textbf{Constraint on interpretation} & \textbf{References} \\
\midrule
\endhead
\bottomrule
\endfoot
Geographically weighted logistic regression & Spatially varying fire-occurrence relationships and future hazard probability & Future scenarios change climate predictors while holding other mapped predictors fixed & \citep{gerberding_mapping_2025} \\
Generalised additive model (GAM) & Nonlinear ignition probability from static and time-varying predictors & Spatial and temporal validation is needed before operational transfer & \citep{moreno_1_nodate} \\
Maxent and logistic GLM & Daily fire-weather suitability, fitted separately by ignition regime & Local relationships may change with season, region and human activity & \citep{de_angelis_modelling_2015} \\
Random Forest & Separate winter and summer susceptibility maps for Italy & Static susceptibility does not describe event intensity or consequences & \citep{trucchia_defining_2022} \\
Canadian Fire Weather Index & Weather-based moisture codes and indices of potential fire behaviour & No explicit topographic term in the index; input weather and danger classes require local evaluation & \citep{arpaci_selecting_2013, zacharakis_environmental_2023} \\
Austrian integrated danger system & Weighted combination of weather, ignition, fuel and terrain layers & Some variants improved ignition prediction; burnt-area correlation remained low & \citep{muller_towards_2020} \\
RISICO & Hourly dead fuel moisture, spread rate and intensity on a 1\,km grid & Fixed live fuel moisture; the spread/intensity analysis used fires above 70\,ha, limiting validation for small fires & \citep{perello_fuelaware_2025} \\
FlamMap 6.2 & Semi-empirical fire behaviour and minimum-travel-time spread & Assumes fixed weather during the simulation; Alpine tests depend strongly on wind and fuel inputs & \citep{malanchini_modeling_2025} \\
PROPAGATOR & Stochastic cellular automaton with time-varying weather & Flexible operational use; intensity and spread-rate outputs differed from FlamMap in three Alpine cases & \citep{malanchini_modeling_2025} \\
LandClim & Landscape vegetation dynamics with fire feedbacks & Long-term vegetation simulations do not provide event-scale operational forecasts & \citep{schumacher_relative_2006, colombaroli_species_2010} \\
LPJmL with SPITFIRE and BASE & European burnt-area scenarios with climate and management assumptions & Continental scenario results require local evaluation & \citep{billing_future_2026} \\
BASE & Empirical European burnt-area model & Reported performance is weaker in mountainous and temperate regions & \citep{forrest_understanding_2024} \\
Debris-flow likelihood model & Logistic probability from severity, terrain, soil and short-duration rainfall & Transfer success in one Mediterranean application does not establish Alpine calibration & \citep{diakakis_exploring_2023, rengers_evaluating_2024} \\
Suppression Difficulty Index & Terrain and fire-behaviour indicators of initial-attack difficulty & Associations validated elsewhere need testing under Alpine suppression practices & \citep{cardil_assessing_2025} \\
Rockfor.net; Root Bundle Model with slideforNET & Rockfall protection and shallow-landslide probability & Inventory-dependent estimates; fire severity and time since fire affect the protection trajectory & \citep{maringer_temporal_2016, gehring_shallow_2019} \\
dNBR, RdNBR, RBR; Composite Burn Index & Spectral change and field assessment of burn severity & Phenology affects retrieval; severity alone may not represent soil hydrological response & \citep{morresi_mapping_2022, negri_wildfire_2026} \\
\end{longtable}
\endgroup

How large the recalibration problem is has been measured directly in an Alpine country. Twenty-two fire weather indices and two raw meteorological variables were tested against 1{,}331 Austrian fire records from 1993 to 2007 at 26 stations spanning 22 ecoregions from colline to alpine zones, and no single index performed best across regions or seasons. In summer, mean daily temperature and the Canadian Build-Up Index were selected most often, with the Build-Up Index dominating once the analysis was restricted to fires above 1{,}000\,m$^2$; in winter a central European index of a different lineage, M68dwd, was selected at 14 of 25 stations for all fires and at 17 for fires above 100\,m$^2$ \citep{arpaci_selecting_2013}. Index choice in this range is therefore conditional on season and on ecoregion simultaneously, which is a stronger statement than the usual observation that indices need local class boundaries. The same analysis reports that 42\% of the fire records lacked a size, which bears directly on the observational problem of Section~\ref{sec:data}.

Alpine practice has answered that conditionality in two ways, and the contrast between them is instructive. Canton Ticino recalibrated locally, and the niche-modelling approach developed there found that meteorological variables alone matched or outperformed the full Fire Weather Index, that the resulting index was not transferable to other regions, and that model selection must be revisited when legislation or socio-economic conditions change the pyrological environment. The same work established that usable indices can be built from small numbers of fires, which matters in regions with sparse records \citep{de_angelis_modelling_2015}.

The Austrian Integrated Forest Fire Danger Assessment System combines five-class weather, human-ignition, lightning, fuel and terrain layers through weighted overlays. Several weighting variants were tested; some improved ignition discrimination and false-alert performance relative to conventional indices, while burnt-area correlations remained low \citep{muller_towards_2020}. Its evaluation supports improved occurrence prediction in that setting, without demonstrating comparable skill for fire size.

These studies identify complementary needs: locally representative weather, suitable fuel descriptions and validation against the intended outcome. Terrain is already represented in some integrated systems, but its influence on local winds can remain unresolved \citep{perello_fuelaware_2025, ozaki_impact_2022}. Austrian fuel measurements also show why imported fuel models require evaluation against managed local forests \citep{neumann_fuel_2022}. Adding predictors alone does not resolve these limitations.

A route out of static fuel representation exists in the Alpine literature, but it has not been carried forward. Landscape simulation for the Swiss Alps using the dynamic landscape vegetation model LandClim treated fire as an emergent property that fed back on vegetation \citep{schumacher_relative_2006, colombaroli_species_2010}, and palaeoecological reconstruction has been combined with dynamic simulation to constrain species response to fire and temperature \citep{colombaroli_species_2010}; more recently, forest-dynamics models coupled to hazard simulation have quantified protective-function loss under climate scenarios \citep{hillebrand_contrasting_2023, moos_climate_2021}. Reconnecting that tradition to fire hazard assessment is, on the evidence assembled here, the highest-value methodological opportunity in Alpine fire science.

\subsection{Data, records and Earth observation}
\label{sec:data}

The records on which Alpine fire science depends differ in resolution by orders of magnitude, and they are best read from the coarsest to the finest. Continental burnt-area products derived from coarse optical sensors resolve at roughly 9\,km and underpin European fire modelling \citep{forrest_understanding_2024}. The European Forest Fire Information System maps individual perimeters and is the source of harmonised European statistics, but its mapping is effective from about 30\,ha upward \citep{muller_forest_2020, beloiu_climatedriven_2026}. Medium-resolution optical archives support disturbance mapping at stand scale, and have been used to classify more than 3.8 million fire, wind and bark beetle events across European forests since 1985 \citep{miguel_unravelling_2026}. National fire databases record events far below any satellite threshold, including a national Austrian register holding several thousand historical and current forest fires \citep{papathomakohle_towards_2025}.

The mismatch between these records is the central methodological fact of Alpine fire science. Mapped fires account for approximately 95\% of the burnt area that European countries report, while representing only a fraction of the number of fires \citep{jrc_advance_2026}. The same gap has been quantified in another European mountain setting, and the ratio is of the same order. Building a danger model for Fennoscandia from 444{,}030 monthly grid observations containing 1{,}439 satellite-detected fire occurrences, that is 0.3\% of the record, the authors found that Norwegian national fire records held about 800 fire data points for 2017 and 2018 where the satellite product held 24, and that a model trained on the national data performed considerably better \citep{bakke_datadriven_2023}. Satellite products systematically underestimate burnt area wherever small fires dominate, and small fires dominate in mountains. For an Alpine country that ratio produces a striking asymmetry: four fires were mapped in Austria in 2025 against a national record running to the low hundreds annually \citep{jrc_advance_2026, papathomakohle_towards_2025}. Any statement about the European fire regime derived from satellite mapping is therefore a statement about area, and any statement about the Alpine regime derived from the same source is a statement about its largest events only.

Earth observation also supplies the severity products on which post-fire assessment depends, and this is where the Alpine case departs furthest from standard practice. Spectral severity indices are the standard spatial proxy for post-fire soil and hydrological condition, and across 113 large Italian fires remotely estimated severity was successfully related to pre-fire land cover, elevation and access \citep{malandra_burn_2022}. But the proxy fails at the point where it matters most for mountains: three years after an Alpine mixed-severity fire the strongest soil water repellency occurred in the low-severity plots, so the spectral signal and the hydrological condition it is used to represent were decoupled \citep{negri_wildfire_2026}. The consequence is that the geomorphic cascade of Section~\ref{sec:water} cannot currently be initialised from the product that operational practice would reach for.

Terrain narrows the Alpine observational window further, and the size of the resulting bias has been measured in the Alps. Steep, deeply incised relief produces illumination contrasts and cast shadow that vary with season and aspect, seasonal snow and persistent cloud occupy part of the usable window, and a difference computed between two dates therefore confounds the change caused by fire with the change caused by solar geometry and by plant phenology \citep{morresi_mapping_2022}. Across ten fires in the mountains of Piedmont following the drought of autumn 2017, severity maps built from phenologically matched reflectance composites rather than from a single pair of cloud-free images raised overall classification accuracy by 4.4 to 5.2 percentage points, reaching 83.7\% at best. They also increased the extent classified as high severity by roughly 19\%, because suboptimal seasonal timing reduces the contrast between healthy and burnt surfaces and so understates what burned \citep{morresi_mapping_2022}. The implication is narrow but consequential: an Alpine severity map produced without controlling for seasonal timing understates high severity by an amount large enough to matter wherever severity class is used as an input rather than as a description.

National registers also hold variables that no satellite product contains, including the number of action forces, helicopters and brigades committed to each event \citep{muller_forest_2020}. They are therefore the only source from which Alpine suppression effectiveness could be estimated, and they are not currently used for that purpose. Detection technology inherits the terrain problem in its own form: automated camera systems that measurably reduce burnt area where they operate were designed for towers overlooking flat ground, and their transfer to complex relief is unresolved, while drones are limited both by sensitivity to wind and by dense canopy \citep{muller_forest_2020}.

Validation practice in Alpine and adjacent fire modelling has improved, and is now reasonably rigorous by the standards of the field. Spatial cross-validation together with an out-of-sample operational test against a severe subsequent season has been applied to national susceptibility mapping \citep{trucchia_defining_2022}; multi-perspective validation explicitly accounting for spatial and temporal structure has been applied to Alpine ignition modelling \citep{moreno_1_nodate}; and negative results have been reported alongside positive ones, as in the finding that integration improves ignition prediction but not burnt-area correlation \citep{muller_towards_2020}. The residual problem is not rigour but statistical power: in a region with few large fires, the events that dominate risk are the events on which models can least be tested.

Seasonal stratification is one of the few design choices that demonstrably helps. Because the Alpine regime is bimodal, models fitted to annual aggregates average over two processes with different drivers. National susceptibility mapping for Italy accordingly produced separate winter and summer surfaces \citep{trucchia_defining_2022}, and analysis of mountain wildfire in south-west China found that different drivers dominate in fire and non-fire seasons and that including the seasonal dimension improved understanding of lagged effects between drivers and occurrence \citep{wang_seasonal_2023}. The same logic applies to the Alps with greater force, since the seasonal contrast there is not merely between active and quiet periods but between two distinct ignition and fuel systems \citep{muller_forest_2020, bercak_combination_2024}.

Deeper still, however, lies the observational threshold. Its magnitude has been measured directly, if not in the Alps. For the Netherlands, the European reference dataset captured 0.2\% of recorded wildfires and 45\% of burnt area \citep{stoof_temperate_2024}; there is no reason to expect better performance over an Alpine fire population that is smaller still and observed through terrain that obstructs detection. The Alpine fire-size distribution described in Section~\ref{sec:hazard} sits almost entirely below the thresholds at which the defining analyses operate: 30\,ha for burnt-area detection \citep{beloiu_climatedriven_2026} or 1\,ha for ignition records \citep{gelabert_assessing_2025}, and the continental burnt-area products underpinning European fire modelling resolve at roughly 9\,km \citep{forrest_understanding_2024}. The Alpine regime therefore sits substantially below the resolution of the instruments used to characterise it, with consequences that run in opposite directions and are easily confused. Detected trends in large Alpine fires are trends in the tail of the distribution and may not represent the regime; and conversely, an absence of detected change may simply reflect that most of the regime is invisible to the detector. Both possibilities argue for the same remedy, which is that Alpine regime detection requires purpose-built records at sub-hectare resolution rather than the reuse of continental products.

The reviewed models represent different parts of the fire process. Weather indices estimate meteorological danger, occurrence models estimate ignition or susceptibility, and spread models describe conditional behaviour. These outputs cannot substitute for one another or for expected loss \citep{muller_towards_2020, trucchia_defining_2022, malanchini_modeling_2025}. Unrepresented fuel change, management and local weather introduce uncertainty, but do not establish that all estimates are lower bounds. The direction of bias depends on the model, setting and omitted process.

The obvious remedy would be to widen them, but the Alpine record cannot support it. Most Alpine fires fall below the observational thresholds discussed above, and the harmonised regional record is short \citep{muller_forest_2020}. The analysis that establishes upslope migration excludes fires below 30\,ha, while conceding that most historical Alpine fires fall below that threshold \citep{beloiu_climatedriven_2026}. Post-fire geohydrological observation amounts to a first inventory with no long-term monitoring \citep{melzner_wildfire-induced_2025}, documented post-fire debris flows in the western Italian Alps number two \citep{tiranti_wildfires_2021}, and systematic Alpine spread-model evaluation rests on three events \citep{malanchini_modeling_2025}. The attribution study of the record 2025 European fire season excludes the Alps and central Europe altogether \citep{keeping_attribution_2026}, as does the annual global assessment for the same season, whose three focal regions are Canadian, Iberian and Chilean \citep{burton_state_2026}. Adding parameters to fit a record this thin will produce models that describe the observations they were fitted to and little else, and that constraint, more than any limitation of method, defines what an Alpine modelling programme can realistically attempt.

A way through is nonetheless visible in the Alpine literature itself. It consists in widening scope through better-chosen structure rather than through more free parameters. Two geometric predictors computable from a digital elevation model and a forest map account for most of the variation in Swiss large-fire statistics \citep{conedera_linking_2024}. A parsimonious additive model with a representative sampling design and threshold-based output achieves high skill on Alpine ignition \citep{moreno_1_nodate}. A locally fitted, low-dimensional index outperforms a transferred complex one \citep{de_angelis_modelling_2015}. Even the post-fire debris-flow problem has been tackled operationally by modifying catchment runoff parameters inside an existing warning system rather than by building a coupled model \citep{tiranti_wildfires_2021}.

Transferability needs to be evaluated for each model and outcome. An empirical post-fire debris-flow model developed in the western United States achieved an area under the ROC curve of 0.913 in one Mediterranean application using European soil data \citep{diakakis_exploring_2023}. A suppression-difficulty framework has also been evaluated across several regions \citep{cardil_assessing_2025}. These are useful precedents for Alpine testing. Neither physical plausibility nor success elsewhere removes the need to assess local calibration and uncertainty.

\subsection{What differs in mountains}

The comparisons in this and the following sections concern the conditions represented by the reviewed studies. Mediterranean and lowland fire systems are themselves diverse, and several share processes found in mountains. We therefore distinguish terrain-related mechanisms from features of the particular Alpine regime.

Within that regime, ignition source and season need to be considered together. Human ignitions dominate annual totals, while lightning becomes important at higher elevations during summer \citep{vacchiano_modeling_2018, conedera_lightninginduced_2006}. A single annual relationship can obscure these differences. Spring peaks also occur in the Carpathians, where present susceptibility is greater at lower elevations \citep{manczinger_fire_2025}. That spatial pattern is not directly comparable with the upward trend over time reported for large fires across European mountains \citep{beloiu_climatedriven_2026}.

Once ignition occurs, slope and terrain-modified winds affect spread. Experiments and simulations outside the Alps show interactions that simple, separate wind and slope corrections may not capture \citep{hayajneh_wind_2026, ozaki_impact_2022}. In three Alpine fires, different spread models produced similar propagation patterns but substantially different intensity and spread-rate estimates \citep{malanchini_modeling_2025}. Thus, agreement on perimeter alone gives limited confidence in the quantities used to estimate damage.

Model evaluation also depends on which fires are recorded. Alpine ecoregions require different danger indices by season, while the many small fires are poorly represented in coarse satellite products \citep{arpaci_selecting_2013, muller_forest_2020}. The same detection problem occurs in Fennoscandia and temperate lowlands, so it is not exclusively mountainous \citep{bakke_datadriven_2023, stoof_temperate_2024}. In mountain terrain, variable illumination and phenology add uncertainty to severity mapping \citep{morresi_mapping_2022}. National event records and locally evaluated satellite products are therefore complementary sources of evidence.

% =======================================================
\section{Climate, Weather and Projections}
\label{sec:climate}

\subsection{Climatic controls on Alpine fire activity}

Alpine fire climate is organised along elevation, along the divide between the moist northern and insubric southern slopes, and along the contrast between dry continental inner valleys and wetter outer ranges. Elevation currently supplies the hard constraint, through snow cover and through temperature-limited fuel drying, and it is precisely this constraint that warming is relaxing. Mountain warming is elevation-dependent, with the snow- and ice-albedo feedback among the principal mechanisms, in an amplification analogous to that of the Arctic \citep{rangwala_climate_2012}; and reduced snow-cover duration is already reshaping the seasonal economy of high-altitude Alpine municipalities \citep{pede_mountain_2022}.

Fuel moisture rather than temperature appears to be the proximate mechanism linking climate to fire. Across European climate regions, multi-month drought indices, soil moisture deficit and vapour pressure deficit are the strongest and most consistent drivers of fire activity, and their interactions are region- and season-specific: productivity deficits matter in southern Europe in summer, whereas in northern Europe summer preconditions combine moisture deficit with productivity \emph{surpluses}, so that fuel production and fuel drying compound rather than trade off \citep{miller_compounding_2025}. The same variable dominates the ignition side. In a five-site European ignition model, the annual anomaly in dead fine fuel moisture content ranked as the most influential predictor at every site \citep{gelabert_assessing_2025}, and vapour pressure deficit is likewise the strongest predictor of burnt area in the mountain upslope analysis \citep{beloiu_climatedriven_2026}.

Drought is not, however, a single quantity, and extreme European fire responds to two structurally different forms of it, a distinction that matters for the Alps because the two Alpine seasons plausibly respond to different ones. Analysis of extreme events in south-eastern France separates a press drought, in which a long-lasting water deficit produced by reduced precipitation and elevated evapotranspiration intensifies wind-driven fires, from a hot drought in which a comparable deficit combines with a summer heatwave to intensify heat-induced fires \citep{ruffault_extreme_2018}. Alpine spring fires, carried by cured herbaceous fuel under foehn, correspond mechanistically to the first; the emerging high-elevation summer fires, associated with extreme atmospheric aridity above 1400\,m \citep{beloiu_climatedriven_2026}, correspond to the second. A danger system calibrated on one will not represent the other, and an Alpine climate projection that reports a single aggregate trend conceals which of the two is changing.

Alpine-specific trend evidence comes from 25 long-term stations distributed across six Alpine countries and grouped into climatically homogeneous regions. Mean annual fire danger, taken as the median of seven danger indices, rose clearly and significantly south and west of the Alps, while north of the Alps and in the inner Alpine valleys a high proportion of trends were not statistically significant; at the 95th percentile, however, almost every station showed a trend towards higher index values \citep{wastl_recent_2012}. Warming is therefore lifting the extreme tail of Alpine fire danger even where the mean is not moving detectably, which is the pattern most likely to be missed by monitoring designed around averages.

Rising danger is not the same as a changing fire regime, and the Alps have much evidence for the first and little for the second. European fire danger showed no clear trend between 1960 and 1999 but a significant increase in southern and eastern Europe from 1980 onward, with correlations between fire-season weather indices and burnt area of order 0.6 \citep{venalainen_temporal_2014}. Detecting a corresponding regime change in the Alps is harder: records are short and administratively fragmented, fires are small enough that recording thresholds dominate the statistics, and the strong seasonal and topographic structure allows compensating changes to cancel in regional aggregates. Under these conditions the absence of a signal is weak evidence of its absence.

A long view from the driest Alpine valleys supports this reading. Reconstruction of the fire history of Valais, a dry continental valley of the Swiss Alps, treats the fire regime as jointly determined by climatic, vegetation and anthropogenic factors and traces regime shifts across the twentieth century in both Valais and Ticino \citep{zumbrunnen_reconstructing_2010}. When some 900 of those fires are related to human and climatic proxies, the controls are seen to change with the socio-economic regime. Between 1904 and 1955 population density was the strongest predictor, followed by wood harvest, and both forest pasturing and wood harvest were associated with fewer fires. Between 1956 and 2006, after the post-war turn from an agricultural to an industrial and service economy, temperature became the leading predictor, which the authors attribute to fire activity concentrating at lower elevations as farming and forestry withdrew from middle and high ground \citep{zumbrunnen_human_2012}. Both models explain only about a third of the variance, so the change is better read as a shift of emphasis than as one driver replacing another. Alpine fire regimes have shifted before under human forcing, and even the emergence of a temperature signal in the Valais record was mediated by land use, which is a caution against attributing any contemporary shift to climate by default.

\subsection{Mountain meteorology: wind systems, foehn and the thermal belt}

At the largest scale, synoptic control on Alpine fire danger has been characterised directly. Combining a sixty-year weather-type classification with station data, danger indices and observed fires across the Alps, elevated danger occurred predominantly under persistent anticyclonic situations or in the lee of the range under katabatic foehn, with the southern Alpine region the most fire-prone; differences between weather types narrowed for indices representing deeper soil layers, reflecting longer drying response times \citep{wastl_largescale_2013}. Foehn is therefore not an incidental local phenomenon but one of the two synoptic configurations that generate Alpine fire danger.

Below the synoptic scale, Alpine wind systems are, in our reading, the mechanism converting small ignitions into fast fires. Foehn can strip tens of percentage points of relative humidity within hours and is associated with the late-winter and spring fires of inner and southern valleys, thermally driven diurnal circulations impose through their anabatic day and katabatic night reversal a spread-direction cycle with no lowland analogue, and relief channels the synoptic flow along valley axes and through cols. Steep Alpine topography exerts a controlling influence on each of these regimes, and thereby on fire type, intensity and burn severity \citep{melzner_wildfire-induced_2025}.

Inversions add a less obvious structure to the same picture. The thermal belt is a recognised feature of the Alpine fire calendar rather than an oddity of one event: in foehn regions north and south of the range, winter fires occur because temperature inversion places a sunny, dry belt above the fog line \citep{muller_forest_2020}. The consequence for danger rating has been measured in the Alps. During the autumn 2011 drought in the Bavarian Alps, fire danger and fire incidents concentrated at mid-slope elevations, in the thermal belt, while valley stations registered only moderate danger; under temperature inversion the diurnal cycles of the controlling meteorological variables differ between valley floor and slope, and once-daily station observation cannot represent that difference \citep{schunk_forest_2013}. The indicated remedy (additional high-elevation stations, or high-resolution meteorological models feeding indices that accept at least hourly input) has not been generally implemented across the Alpine danger-rating systems reviewed in Section~\ref{sec:hazard}.

\subsection{Lightning, convection and the summer ignition mode}

Lightning frequency was among the strongest predictors in Alpine-scale hazard mapping, alongside mean annual temperature and precipitation seasonality \citep{gerberding_mapping_2025}. Its importance is consistent with the summer ignition regime, although predictor importance in one model does not establish the absence of similar relationships elsewhere.

Multi-decadal records have characterised the lightning regime across Valais, the Grisons, Valle d'Aosta and Ticino. Lightning fires there occur from May to October, with 78\% of events in July and August. They ignite at higher elevation on steeper relief than anthropogenic fires, and are harder to extinguish because they burn underground. In hot dry summers such as 2003 they became markedly more frequent, with significantly greater burnt area and higher suppression costs \citep{conedera_lightninginduced_2006}. Lightning ignition is itself terrain-structured, through elevation-dependent cloud-to-ground discharge frequency, altitudinal fuel belts and strong vertical gradients in precipitation \citep{baranovskiy_deterministic_2021}. This is reflected in Alpine operational practice, where lightning fire hazard is classified jointly by elevation band, slope, aspect and vegetation type rather than by strike density alone \citep{muller_towards_2020}. Whether ignition potential and fuel receptiveness will increase together is the compounding question: rising extreme fire weather alongside increased thunderstorm activity has been projected for the Dinaric setting, where lightning-induced fires are rare but extremely costly to electricity transmission networks \citep{kyriakopoulos_impact_2023}.

\subsection{Climate projections and what they omit}

Alpine-scale hazard mapping used geographically weighted binary logistic regression to relate fire occurrence to environmental and human predictors \citep{gerberding_mapping_2025}. Future maps changed the climate inputs under SSP scenarios while holding the other mapped predictors fixed. The projected increases reported in Section~\ref{sec:intro} therefore describe conditional changes in occurrence probability, rather than simulated fire spread or damage. Fire-weather projections answer a related but different question by estimating the frequency of dangerous meteorological conditions \citep{miller_climate_2024}.

Both approaches leave parts of the evolving fire system outside the calculation. The Alpine occurrence scenarios do not dynamically simulate disturbance-driven fuel change or future land management \citep{gerberding_mapping_2025}. Continental burnt-area modelling addresses some broader controls, but BASE reports weaker performance in mountainous and temperate regions \citep{forrest_understanding_2024}. These differences in model purpose and spatial support need to remain explicit when interpreting projections.

Anticipation faces a further limit, although the evidence for it is indirect. A global analysis of thirteen climate teleconnections found that collectively they modulate around half of global burnt area, but that these relationships are markedly weaker in Europe than on other continents, with only the North Atlantic Oscillation and the East Atlantic pattern showing consistent associations in temperate regions \citep{cardil_climate_2023}. Seasonal fire forecasting from remote climate modes will therefore deliver less in the Alps than in regions where a single mode dominates, which places correspondingly more weight on land-surface and fuel-state monitoring.

\subsection{What differs in mountains}

Regional warming affects fire weather across both mountain and lowland Europe, but relief changes how regional conditions are expressed locally. Alpine station analyses distinguish northern, southern and inner-valley fire climates, with trends that differ between mean and extreme danger \citep{wastl_recent_2012, wastl_largescale_2013}. Regional averages can therefore conceal locally important changes.

Within those climatic regions, foehn, inversions and valley winds alter fuel drying and the wind experienced by a fire. During one Bavarian Alpine drought, elevated danger at mid-slope was poorly represented by valley stations \citep{schunk_forest_2013}. Evidence from Tasmania independently shows how ruggedness increases the importance of vertical atmospheric structure \citep{ozaki_impact_2022}. The implication is to evaluate interpolated weather against conditions at exposed slopes, rather than to assume that interpolation is uniformly inadequate.

At higher elevations, snow-cover duration additionally constrains the period available for fuel drying. Warming and vegetation infilling may change this constraint together, although their combined effect on ignition requires local measurement \citep{rangwala_climate_2012, frei_evidence_2023}. Snow also affects northern lowlands; the mountain distinction is its strong variation with elevation and aspect. Seasonal lightning ignitions add a further spatial dependence on terrain and vegetation \citep{conedera_lightninginduced_2006, muller_towards_2020}.

These local controls limit how directly a regional projection can inform site management. Alpine hazard projections indicate increasing climatic suitability for fire, but do not resolve all changes in fuels, ignitions or management \citep{gerberding_mapping_2025}. We interpret them as conditional scenarios whose local consequences require further assessment.

% =======================================================
\section{Water, Drought and Water Use}
\label{sec:water}

\subsection{Drought, fuel moisture and the physiological chain}

Drought reaches fire through a chain of separable links rather than as a single driver, and the fire literature has generally attended to the link closest to fire while the forest health literature has attended to the rest.

That closest link is fuel moisture, which is well established though not well modelled. A global evaluation across 1{,}603 measurements from 153 experiments in 21 countries found vapour-pressure-deficit models performing competitively over the full moisture range, with mean absolute errors near 4\% and coefficients of determination around 0.4, and machine-learning alternatives improving on them only marginally; below 10\% moisture, however, every approach degraded to a coefficient of determination of about 0.11, which is precisely the range in which spread responds exponentially to dryness \citep{rodrigues_vpdbased_2024}. Fuel moisture is predicted worst where it matters most. Within that limitation, dead fine fuel moisture anomaly is the strongest single predictor of ignition likelihood across European pilot regions \citep{gelabert_assessing_2025}, and multi-month drought indices together with vapour pressure deficit dominate burnt area \citep{miller_compounding_2025, beloiu_climatedriven_2026}.

Further back in the chain, drought acts physiologically and converts living trees into standing dead fuel. In Scots pine at the dry margins of its range, which includes the inner Alpine valleys, drought-induced mortality proceeds through carbon starvation, as reduced assimilation under water stress depletes the reserves that also fund defence, leaving trees susceptible to pest and pathogen attack \citep{giuggiola_droughtinduced_2010}. Experimental work isolates the interaction: in Norway spruce, simultaneous drought and root-rot infection create competing demands on carbon, since compartmentalising the pathogen and maintaining turgor draw on the same pool, while the infection itself reduces functional sapwood and so compounds water stress; needle starch concentration predicted which saplings died \citep{gomezgallego_interaction_2022}. Scots pine showed no such interaction in the same experiment, which matters because the two species occupy different Alpine positions and would therefore be expected to fail by different routes.

Whether such decline proceeds gradually or abruptly is not settled, although the comparative evidence favours abruptness. In central Chile, more than 90\% of a Mediterranean-type forest browned within fewer than a hundred days under conditions unprecedented across all five Mediterranean-type ecosystems, with long-term recovery uncertain given a continuing decline in regional water balance \citep{miranda_widespread_2023}. The setting is not Alpine and the magnitude does not transfer, but the mode does: drought-driven decline can be synchronous and regional rather than incremental, which is the behaviour that would convert a gradual fuel-loading trend into a step change.

The chain is completed when dead and weakened trees become fuel, a link that the Alpine hazard literature asserts but its models leave out: projections of Alpine fire hazard note that drought, heatwaves and storms favour bark beetle outbreak and windthrow and thereby increase the availability of flammable material, while conceding that the pathway is absent from the projection \citep{gerberding_mapping_2025}. Since disturbance is constitutive of Alpine forest dynamics rather than exceptional \citep{bebi_changes_2017}, and since European forest disturbance frequency has been increasing over three decades while severity has been falling \citep{senf_mapping_2021}, this omission is unlikely to be small.

Because the same stress drives both mortality and fuel loading, the chain closes on a single silvicultural lever: in inner Alpine valleys, intervention against drought stress and intervention against fuel are substantially the same operation, an alignment developed in Section~\ref{sec:forest} and one that management planning does not currently exploit.

\subsection{Soil alteration and the limits of severity as a proxy}

Post-fire assessment often uses mapped severity to infer changes in vegetation and soil. Fire can remove litter, deposit ash, alter water repellency and weather exposed rock \citep{santi_wildfire_2022}. Whether a severity class captures the hydrological consequences requires field evaluation; it cannot be assumed from canopy change alone.

At one Alpine Scots pine site, soil sampled three years after a mixed-severity fire retained changes in surface cover, carbon and aggregate properties. These changes were not proportional to mapped severity, and water repellency was strongest mainly in low-severity plots with abundant uncharred organic matter \citep{negri_wildfire_2026}. This case cautions against using severity alone to estimate soil hydrological condition. It does not invalidate severity mapping for all post-fire applications.

The finding cuts against the standard post-fire workflow, in which a severity map serves as the spatial proxy for hydrological and geomorphic change. It also complicates a common assumption in fuel management, since it implies that a low-intensity burn is not automatically hydrologically benign.

That so basic an inference has gone largely untested reflects a wider neglect of mountain soils. A comparative review of the soils of the Australian Alps records that above-ground ecosystems there received far more attention than the soils beneath them, even though high-elevation soil types are disproportionately important to hydrology and slope stability and disproportionately vulnerable \citep{wilson_distribution_2022}. The European Alpine literature shows the same asymmetry, and since the Alpine cascade runs through soil condition rather than canopy loss alone, the imbalance has consequences beyond the soil science itself.

Alpine natural-hazard consequences of fire were catalogued before the recent inventories, and the catalogue is broader than debris flow alone: new avalanche-prone slopes, mudslides, increased runoff, rockfall and soil erosion, with steep subalpine slopes carrying a thin humus layer over limestone identified as the setting where complete soil loss and failure of regeneration are most likely \citep{muller_forest_2020}. What the subsequent decade added was quantification rather than recognition.

\subsection{Erosion, debris flow and sub-threshold triggering}

Post-fire debris flow is a potentially damaging pathway in Alpine valleys. Two documented events in the Western Italian Alps, Pallanzeno in 2005 and Bussoleno in 2018, followed moderate, localised rainfall over recently burnt catchments \citep{tiranti_wildfires_2021}. Their deposits contained fine particles and burnt wood, consistent with fire-altered surface material being redistributed into channels. The small number of cases limits regional estimates of frequency or magnitude.

One of those two events, Bussoleno, has since been analysed numerically, and it supplies parameters that make the mechanism quantitative even though the sample remains a single catchment. An October 2017 fire burnt over 95\% of a 1.37\,km$^2$ catchment; after a spring 40\% wetter than normal, a rainfall of 12\,mm in 30 minutes (an intensity of 24\,mm\,h$^{-1}$, a five-year return period) mobilised roughly 20{,}000\,m$^3$ of sediment and woody debris some eight months after the fire, destroying two buildings and flooding around 150 residents from the alluvial fan. The mechanical shear strength of burnt soils was reduced by a factor of three, and the triggering intensity lay below the threshold of more than 30\,mm\,h$^{-1}$ typical of unburnt basins of that type \citep{laporta_postwildfire_2025}. Fire does not simply increase the response to a given rainfall; it lowers the rainfall required, which moves the hazard into a return period that occurs several times per decade.

Monitoring after the Grizzly Creek Fire in Colorado provides a useful methodological comparison. The first-year warning threshold was 25.9\,mm\,h$^{-1}$, and 89\% of observed debris flows followed rainfall above it \citep{rengers_evaluating_2024}. However, this threshold used a 15-minute duration, whereas the Bussoleno intensity above averages 30 minutes; the numerical similarity is not a direct validation. The second-year threshold of 33.7\,mm\,h$^{-1}$ also used a different probability criterion, changing from 50\% to 75\%. Eight storms exceeded it without producing debris flows as vegetation recovered. The study supports updating post-fire assessments, while showing uncertainty in operational thresholds. Its volume model overestimated measured erosion by a median factor of 4.4.

One sediment source specific to rocky burnt terrain has now been quantified in a European mountain setting. After a fire that burnt 131.64\,km$^2$ of Sardinia at temperatures above 800\,\degree{}C, with 63.1\% of the area at moderate-high or high severity, thermal spalling was recorded at 69 rock exposures, removing between 0.07\% and 74.57\% of the exposed surface with a median of about 20\%, in fragments from 3\,mm to 13\,cm. The material entered the channel system: spall fragments made up between 1.3\% and 36.6\% of the deposits of the 33 debris flows that occurred across 21 sub-basins in the following two years, triggered by three rainstorms, and freshly spalled angular cobbles accounted for 10 to 30\% of all cobbles in the deposit \citep{pala_sediment_2025}. Fire in rocky terrain therefore manufactures transportable sediment directly, rather than only exposing what was already there, and the process was insensitive to the distinction between high and moderate-high severity while being sparse below it.

Sediment supply, rather than runoff alone, appears to set event magnitude, though we note that the demonstration is imported rather than Alpine. Lidar differencing of the 2018 Montecito flows in California showed 550{,}000\,m$^3$ of previously stored channel sediment progressively entrained and partly redeposited, with scour concentrated where coarse colluvium had accumulated and persisting beyond the mountain front, implying that volume is controlled by the reservoir available for scour and hence by catchment history \citep{morell_sediment_2021}. For the Alps this means post-fire debris-flow magnitude is a function of the antecedent state of the channel network as much as of the fire.

\subsection{Water supply, water quality, snow and ice}

The cascade does not stop at the fan apex but continues into the water system that the range supplies. Alpine catchments supply drinking water, hydropower and irrigation to populations far larger than those living in the mountains, and the same processes that mobilise sediment onto a road deliver it, with its associated chemistry, into storage.

How large the water-quality signal is has not, however, been measured in the Alps. It has been modelled for a Mediterranean supply reservoir under a scenario burning around 30\% of its contributing area: median suspended solids rose by a factor of 4.7 and maximum nitrate and phosphate concentrations by a factor of four, although dilution within the reservoir removed roughly half the nitrate and nine tenths of the phosphate before the dam wall, so that drinking-water standards were not breached \citep{basso_potential_2021}. The mobilisation and the dependence of the outcome on storage volume are general enough to transfer to the Alps, but the buffering capacity is contingent, and the same study notes that dilution capacity falls during drought, that is, precisely when the antecedent conditions for fire were present. Systematic review of fire effects on ecosystem services in a southern European setting found hydrological cycle regulation to be the one category in which every recorded entry was negative, with the most pronounced effects in the first two post-fire years \citep{misseyanni_impact_2025}, and water regulation likewise features among the services assessed as at risk from wildfire in Mediterranean forest landscapes \citep{lecinadiaz_assessing_2021}.

Drought can increase dependence on glacier melt at the same time as it raises fire danger. During the 2022--2023 snow drought, melt contributions rose markedly in two Italian Alpine catchments but did not fully compensate for the precipitation deficit \citep{leone_20222023_2026}. Fire-related changes to snow and ice could interact with this dependence. We treat that interaction as an Alpine research question, rather than an established positive feedback.

Snow cover limits the period available for drying at high elevations, while vegetation infilling changes the amount and continuity of potential fuel \citep{rangwala_climate_2012, frei_evidence_2023}. Fire can also affect subsequent snow accumulation and melt. The available evidence for this reverse relationship comes mainly from outside the Alps.

Canopy loss changes snow interception and surface energy exchange, while deposited debris and black carbon can reduce albedo \citep{koshkin_wildfire_2022}. At Athabasca Glacier in the Canadian Rockies, lower ice albedo persisted after wildfire activity and increased modelled melt by slightly more than 10\% in later smoke-free seasons. During smoky periods, reduced incoming radiation partly offset that increase \citep{aubrywake_fire_2022}. Smoke and surface darkening therefore have competing effects that should be represented separately.

A comparison of one burnt and one unburnt basin in the northern Sierra Nevada found increased snow-covered area after fire but earlier basin-average melt-out \citep{micheletty_application_2014}. This distinguishes snow extent from melt timing; snow-covered area alone does not measure stored water volume. Alpine studies would need to resolve accumulation and energy balance before inferring consequences for seasonal water supply.

\subsection{What differs in mountains}

Post-fire runoff, erosion and water-quality deterioration occur in many landscapes. Steep catchments can additionally connect sediment sources to settlements through debris-flow channels and alluvial fans \citep{santi_wildfire_2022, tiranti_wildfires_2021}. Fire may change both the supply of movable sediment and the rainfall required to mobilise it. This matters for mountain protection because downstream assets can be affected without burning.

The Bussoleno case shows that damaging flow can follow rainfall below the threshold normally associated with comparable unburnt catchments \citep{laporta_postwildfire_2025}. However, one catchment does not establish a regional threshold. Monitoring in Colorado shows that threshold performance changes as vegetation recovers, supporting repeated assessment rather than a fixed post-fire warning value \citep{rengers_evaluating_2024}.

Recovery must also be defined for the process being assessed. Soil hydrological properties, sediment availability, root reinforcement and rockfall interception recover on different timescales \citep{negri_wildfire_2026, gehring_shallow_2019, maringer_temporal_2016}. Persistent canopy loss is relevant to protection, but does not by itself establish that a debris-flow rainfall threshold remains lowered for the same period. We therefore avoid treating these observations as measurements of one common hazard window.

In snow-fed and glaciated catchments, changes in water supply add another consequence pathway. Alpine drought observations demonstrate increased reliance on glacier melt, while studies outside the Alps document fire-related changes in snow cover and ice albedo \citep{leone_20222023_2026, micheletty_application_2014, aubrywake_fire_2022}. The combined Alpine effect remains a research question. It cannot be assumed to increase melt uniformly, because atmospheric smoke can partly offset the radiative effect of darker ice.

% =======================================================
\section{Forest: Change, Ecology and Management}
\label{sec:forest}

\subsection{Forest change, composition and fuel connectivity}

The landscape now burning is a recent artefact of land use, and recognising it as such determines which policy lever can move it. Centuries of grazing, litter collection and fuelwood extraction held Alpine forests open in a mosaic that limited spread irrespective of weather; the retreat of that regime has reassembled the fuel bed \citep{bebi_changes_2017}. Expansion is not uniform. It is fastest where recent agricultural abandonment coincides with climatically favourable conditions for establishment, so that anticipating future fuel continuity requires knowing how the timing of land-use change intersects with growth conditions, not simply where abandonment occurred \citep{garbarino_contrasting_2020}. The generalisation that abandonment monotonically increases fire hazard should nonetheless be treated as a regionally supported finding rather than a law, since for disturbances other than fire the reported effects of abandonment are inconsistent and sometimes opposite in sign \citep{mantero_influence_2020}.

Expansion also has a vertical component. In one intensively surveyed Central Swiss valley, the treeline advanced by an average of 10\,m per decade between the early 1970s and 2012, with local maxima of 42\,m per decade; decisively, only 21\% of new trees represented advance beyond the former treeline while 79\% were infilling below it \citep{frei_evidence_2023}. Fuel is therefore not merely climbing but densifying, and densification below the treeline is what converts scattered stands into continuous fuel. At Alpine scale, 24\% of the 10.5 million hectares of forest is in an open-canopy structural state, concentrated at high elevations, forest edges and warm dry sites \citep{stritih_alternative_2023}, the same positions where fuel is now filling in.

Where that regrowth happens is, in turn, economically determined and locally variable. A spatially explicit analysis of the Swiss mountains found forest re-growth concentrated where cultivation costs were high and yield potential low, with non-linear relationships to slope steepness and soil stoniness, and concluded that general policy measures applied across the whole mountain area are unsuitable for preventing abandonment and re-growth, which instead require attention to local characteristics \citep{gellrich_agricultural_2007}. Swiss forest area has risen by roughly a third over the last 150 years, most of it by natural re-growth rather than planting. Fuel accumulation in the Alps is therefore the spatial signature of agricultural marginality, which is also why it concentrates on the steep ground where fire runs longest.

That the present composition of a European mountain forest can be a land-use artefact rather than a climatic equilibrium has been demonstrated on a multi-millennial record. At 1{,}390\,m in the Northern Apennines, species-rich mixed silver fir forests dominated from about 10{,}500 to 5{,}500 years ago, with fir reaching up to half the pollen assemblage, before beech expanded and eventually replaced them entirely. Charcoal influx and temperature each explained about a tenth of the variation in the vegetation record. Fir, elm, ash and lime all responded strongly negatively to burning, while beech tolerated light to moderate fire and declined under heavy fire \citep{moralesmolino_longterm_2021}. The conclusion drawn there is that today's monospecific beech forests are the product of multi-millennial land-use intensification superimposed on late Holocene cooling, and that the displaced mixed fir forests were the more drought-tolerant system. The parallel with the Alpine fuel bed is structural rather than incidental: in both ranges the stand that now carries the fire is an artefact of how people used the land, which is also why the correction is a land-use correction.

Whatever its origin, the reassembled fuel bed has a direct behavioural consequence. Where fuel connectivity and slope length jointly govern large-fire likelihood \citep{conedera_linking_2024}, expansion onto steep marginal ground that was abandoned first increases exactly the quantity that matters. This is the mechanism by which a landscape can become more fire-prone without any change in weather.

How much fuel that bed holds has been measured destructively in one Alpine country. Across 93 plots in Austrian Norway spruce, Scots pine and black pine forests, fuel load depended on region and forest type rather than on stand age, dead surface fuel increased and live understorey fuel decreased with basal area, and spruce stands carried the highest loads \citep{neumann_fuel_2022}. The forest type that burns most often in Austria, black pine on the dry eastern lowland margin, carried the lowest load but combined open canopies with greater shrub and grass cover and the deepest fuel bed. The authors take this to indicate that fuel structure, together with the drier eastern climate, contributes to its fire frequency. The measurement leaves out the layers that matter most once the bed dries. The duff layer, which was not sampled, is estimated from the literature and from Austrian forest-floor inventories to add between 15 and 70\,t\,ha$^{-1}$ that can burn and smoulder under drought. Canopy fuel, roughly 10\,t\,ha$^{-1}$ in spruce stands, is consumed only in the crown fires that are currently rare in central Europe \citep{neumann_fuel_2022}. Both are reserves that a drier climate would bring into play.

Fire in the Alps also arrives into a disturbance regime rather than a quiescent forest, and the interactions are structured rather than additive. Wind, snow break, bark beetle and avalanche shape Alpine stands continuously \citep{bebi_changes_2017}, and at Alpine scale a quarter of forest area is already in an open-canopy structural state, with disturbance initiating transitions to that state in roughly half of cases and recovery to closed canopy taking around 35 years in most forest types \citep{stritih_alternative_2023}.

Disturbance regimes now have a measured European baseline resolved by causal agent. Classification of more than 3.8 million fire, wind and bark beetle events across European forests between 1985 and 2023 separates fire regimes on a north--south gradient, severe and rare in boreal and temperate forest against moderate, large and frequent in Mediterranean and transitional forest, and separates wind and beetle disturbance into four regimes of which the frequent and small predominates in Central Europe \citep{miguel_unravelling_2026}. The result that bears hardest on Alpine assessment is not the classification but its instability: about 15\% of grid cells changed regime within the record, fire frequency and severity fell in temperate and Mediterranean forest between the first and second period, and wind and beetle disturbance then rose in severity in temperate forest between the second and third. A disturbance regime is not a fixed property of a region on which a fuel model can be built once.

Disturbance susceptibility appears, moreover, to be partly a management artefact rather than a purely climatic one. The MANFRED programme, a transnational Alpine Space collaboration spanning fifteen partners in six countries, concluded that recorded storm damage in Alpine forests is not driven by climate alone, and that the prevailing silvicultural practice of allowing increasingly dense forests to mature has itself produced stand types that are markedly vulnerable to storms \citep{jandl_management_2013}. The fuel and the fragility have a common origin in how the forest has been managed.

Fire and windthrow can leave different protective structures. At one windthrown Dolomite site, modelling indicated that fallen logs and stumps could provide substantial rockfall barriers \citep{costa_biological_2021}. Fire can consume litter and weaken stems, but does not necessarily remove all biological legacies. Southern Alpine beech research explicitly recommends retaining suitable standing and fallen dead trees for their temporary protective contribution \citep{maringer_temporal_2016}. Post-disturbance protection depends on the severity of damage and the condition and arrangement of surviving material.

It follows that salvage decisions after fire cannot simply import the consensus developed after windthrow. Where that consensus advises against salvage in order to retain protective legacies \citep{costa_biological_2021}, the evidence base is itself less categorical than often assumed: eighteen years after a Tatra windstorm, salvaged areas retained an eighth of the deadwood volume of unsalvaged areas but supported higher regeneration diversity, 5.40 against 3.85 species per plot, with unsalvaged plots dominated by Norway spruce \citep{seben_impact_2024}. The trade-off between protective legacy, regeneration composition and future fuel is real, differs between disturbance types, and has not been evaluated for Alpine post-fire sites.

A comparative case from outside Europe marks the boundary of the Alpine trajectory. In the southern Appalachian Mountains decades of suppression preceded a regional outbreak of more than a dozen large wildfires during one of the warmest and driest droughts of the last century, in a landscape whose historical mosaic had been maintained by frequent low- and mixed-severity fire and whose composition had already been reorganised by land-use change \citep{reilly_drivers_2022}. The Alps never had a frequent-fire regime to suppress, so the trajectory is not the same one; what the comparison establishes is only that a landscape whose fuel structure has ceased to match its disturbance history can support a regional outbreak once drought supplies the trigger.

What grows back matters as much as how much, because species composition is both an input to fire and an outcome of it. Alpine tree species differ sharply in their capacity to survive fire, and the composition of Alpine forests is being reorganised on a timescale comparable to that of regime change. Simulation of Swiss Alpine landscapes indicates that climate-driven vegetation shifts will exert a larger influence on biomass and species distribution than management, with fire emerging as a property of the simulated landscape that then feeds back on vegetation \citep{schumacher_relative_2006}.

Vulnerability is trait-mediated, and it now has independent experimental support. Simulation across 4438 French Alpine forest plots showed cold-season fires, about half of all fires there, having limited impact except on low-elevation southern stands, average summer conditions producing highest mortality in mixed and broadleaved low-elevation forests with up to 75\% mortality in coppice, and extremely dry summers producing high mortality across all communities, with the lowest mortality in thick-barked, high-canopied larch stands of the inner Alps \citep{dupire_vulnerability_2019}. Controlled heating of stem discs from six Central European species reproduces that ranking from first principles: bark thickness is the most robust predictor of the time taken for the cambium to reach the lethal threshold, placing larch, Douglas-fir and Scots pine above silver fir and Norway spruce, with beech lowest, and thermal resistance declining with bark moisture and with height on the stem \citep{mutterer_bark_2026}. Two very different methods converge, which is a stronger basis than either alone.

The relationship also runs in the opposite direction, from fire to composition. After a stand-replacing fire in the Rhone valley, recruitment favoured pioneer birch, aspen and willow over the pre-fire larch, spruce and Scots pine, and seedling density tracked moisture availability rather than fire intensity, leading to the conclusion that drought is the limiting factor for post-fire recruitment in the Central Alps \citep{moser_potential_2010}. Palaeoecological records extend the same logic across millennia, with fire and grazing rather than temperature alone determining whether the timberline attains its climatic potential \citep{colombaroli_species_2010}, and with regionally distinct Holocene fire--vegetation trajectories reconstructed for the Dinaric ranges \citep{cagliero_effect_nodate}.

\subsection{Protective function, biodiversity and the recovery window}

Consequence on an Alpine slope turns on what fire removes, how long it stays removed and what else changes in the meantime. Of the post-fire processes involved, rockfall is the most specific to the Alps and the least represented internationally; as recently as 2019 no work on post-wildfire rockfall activity had been published for the eastern Alps, and even in the better-studied western United States the process had received little attention relative to debris flow \citep{melzner_brief_2019}. The thermal mechanism has a temperature scale attached to it. Review of three Spanish cases shows the change operating simultaneously in the source area, the propagation zone and the affected area. In the source area, mineral expansion, dehydration and decarbonation open microcracks, with granites losing strength appreciably above 500\degree{}C as quartz approaches its phase transition at 573\degree{}C, and limestones altering from about 300\degree{}C. In the propagation zone, loss of vegetation removes the energy absorption that reduces both the reach and the frequency of blocks. In the affected area, the cables and barriers installed to arrest rockfall are themselves degraded by heat \citep{sarro_effects_2021}. The same review records that the restitution and friction coefficients used in rockfall trajectory models have not been standardised for post-fire conditions, so the parameters that would carry a burnt slope through a trajectory model are the parameters least constrained by measurement.

Vegetation enters the problem on two timescales, through the immediate loss of the litter and vegetation that arrest small blocks and through the slower loss of standing stems that intercept falling material, so exposure can rise for years as fire-killed trees decay and fall. Trajectory modelling isolates the vegetation term and shows it to be strongly block-size dependent: small blocks around 0.01\,m$^3$ that mostly came to rest on a vegetated slope travel much further once cover is removed, while for blocks an order of magnitude larger the kinetic energy overwhelms the vegetative resistance \citep{asteriou_assessing_2025}. Protective function is thus lost first for the frequent, small events that dominate long-run exposure.

How fast that function is lost, and for how long, has been quantified for beech on the southern side of the range, and the numbers are more severe than the general disturbance-window literature implies. Across 244 plots in Ticino and Piedmont between 600 and 1{,}700\,m, root reinforcement after high-severity fire fell by roughly 70\% within five years and by 98\% by year fifteen, while moderate severity produced about 15\% loss in the first five years and 60\% by year twenty-three; low-severity burns remained stable. Modelled protective capacity against shallow landslides fell below a quarter of the unburnt value after high-severity fire irrespective of soil and slope class, and the elevated landslide probability persisted for at least forty years \citep{gehring_shallow_2019}. The rockfall function follows a similar trajectory on a different clock: across 34 burnt beech stands in the southern Alps, low-severity burns retained close to unburnt protective capacity, while moderate and high severity reduced it significantly and did so most strongly between ten and thirty years after the fire \citep{maringer_temporal_2016}.

These results show that protective capacity may decline gradually as damaged trees die, before regeneration can replace them. An assessment made immediately after fire can therefore miss a later period of reduced protection. Root reinforcement and rockfall interception also follow different trajectories, so their durations should not be equated with shorter runoff and erosion responses.

Numbers exist for the economic weight of that function at one site. For a protection forest above Martigny in the Valais, coupled forest-dynamics and three-dimensional rockfall simulation places annual rockfall risk at 19{,}100\,CHF with the forest present and 95{,}500\,CHF without it, a factor of five; under climate scenarios the risk rises to 59{,}250 and 70{,}100\,CHF per year within 150 years as composition shifts from beech- to pubescent-oak-dominated, and small blocks reaching the railway move from once per decade to once every three years \citep{moos_climate_2021}. Simulation across an elevation gradient in South Tyrol shows the loss is elevation-contingent: montane forests decline drastically in avalanche protection under severe warming, with drought-induced biomass reductions reaching 97\% in dry inner valleys, while subalpine stands remain stable or improve \citep{hillebrand_contrasting_2023}.

For risk assessment, duration matters alongside the magnitude of the change. Reviews report post-fire geomorphic responses lasting several years and, in some settings, much longer \citep{santi_wildfire_2022}. Alpine forest studies add evidence of delayed and prolonged loss of protective structure, but each process needs its own recovery estimate.

Independent lines of evidence converge on that conclusion. Soil properties remained altered three years after an Alpine mixed-severity fire, including properties that ordinarily recover fast \citep{negri_wildfire_2026}. Recovery of the canopy that restores protective function is drought-limited at exactly the low and mid elevations where assets concentrate, requiring precipitation far above the long-term average for successful establishment \citep{moser_potential_2010}. At Alpine scale, moreover, although 72\% of disturbed forests return to a closed-canopy state within 35 years, submediterranean forests remain open for decades \citep{stritih_alternative_2023}. At the upper margin the window may not close at all, although the observation comes from one massif: on burnt subalpine slopes in the Northern Limestone Alps at 1800--1950\,m, species richness was higher in burnt plots but tree seedlings were nearly absent and no meaningful mountain pine re-establishment occurred over eighty years, a state better described as arrested succession than as delayed recovery \citep{degiuli_arrested_2024}.

The length of that window is therefore partly a management outcome, and current Alpine restoration practice is poorly matched to it (Section~\ref{sec:fuelbreaks}).

Fire in the Alps is not, however, uniformly destructive, and the review literature is unusually candid about the risk of assuming otherwise. A systematic assessment of mountain wildfire against fifty ecosystem services found that current fire regimes mostly reduce service provision, that effects on cultural services are mixed, that the dominant mechanism is change in vegetation composition and structure, and that a bias toward negative characterisation may be built into how the questions are posed \citep{vukomanovic_systematic_2019}.

Where Alpine forest sits in a European ranking of ecological vulnerability can now be estimated. A continental assessment combining ecological value, the resistance of vegetation to fire and post-fire recovery time places temperate conifer forest second among European biomes, behind tundra and ahead of boreal forest, and the term responsible is recovery time rather than resistance \citep{arrogantefunes_mapping_2024}. Non-forested vegetation, which includes the grassland, heathland and shrubland of the subalpine and alpine belts, carries no coping capacity in that framework, a modelling convention rather than a measurement, but one that follows from the absence of the bark and crown traits on which resistance is defined. Taken together the two results locate the Alpine problem in the recovery term: what is lost is not resisted at the moment of the fire and is returned slowly afterwards, so that the integral of lost ecological value can be large where the burnt area is small.

Abundance and composition respond differently, and conflating them yields the wrong conclusion in either direction. Across two inner-Alpine Swiss valleys, high-severity fire roughly halved bird abundance (21.68 against 41.59 territories per 10\,ha) without significantly changing species richness; fourteen species were more abundant in burnt stands and five in unburnt ones, red-listed and priority species were substantially more common in the burnt stands, and the effect strengthened with time since fire, many species peaking eight to twelve years after the event and declining as the vegetation densified \citep{rey_effects_2019}. Fire in these valleys removed individuals and rearranged the community rather than eliminating species, and the species it favoured were disproportionately those of conservation concern. The same asymmetry appears in the vegetation, where post-fire subalpine slopes in the Northern Limestone Alps carried significantly higher vascular plant richness than unburnt controls, including alpine specialists absent from the unburnt plots, even as tree regeneration failed entirely \citep{degiuli_arrested_2024}.

What neither measurement contains is mortality. Individual animals are killed by burning, smoke inhalation and heat stress, and the taxa most exposed are those that cannot outrun a front or that depend on structures the fire removes: small mammals, cavity-nesting birds, reptiles, and aquatic organisms reached through post-fire sediment delivery. Hollows used for denning and nesting are lost with the stems that contain them, and poor dispersers are slowest to recolonise \citep{demiguel_fires_2021}. These mechanisms are general rather than Alpine, and the only order-of-magnitude figures available come from a single non-European event. No Alpine estimate of faunal mortality exists in the literature assembled here, so the animal term in the Alpine loss function is not a small quantity but an unmeasured one.

These findings specify the risk framing of this review rather than softening it, because they show the Alpine loss function to be composite. The protective function of forest on slopes above assets is spatially concentrated, economically quantifiable \citep{moos_climate_2021} and slow to recover \citep{stritih_alternative_2023, degiuli_arrested_2024}. The ecological value integrated over that recovery period is high for the forest type that dominates the Alpine belt \citep{arrogantefunes_mapping_2024}. The biodiversity term is negative for abundance and for closed-forest specialists and positive for open-habitat and disturbance-dependent species \citep{rey_effects_2019, degiuli_arrested_2024}. The faunal term, by contrast, remains unquantified. Collapsing these into one quantity produces the two opposing errors visible in the wider literature: fire as uniformly destructive, and fire as ecologically benign because some species benefit.

\subsection{Managing for protection: silviculture, species and fuel treatment}

Alpine forestry is unusual in that its principal product is often not wood. Across large parts of the range the managed objective is the maintenance of a stand structure capable of arresting rockfall, initiating and anchoring snowpack against avalanche release, and holding soil on steep ground, with timber a secondary output constrained by extraction cost in terrain where harvesting is technically demanding and economically marginal \citep{dinca_timber_2025}.

That the protective role should dominate Alpine forest management is long established: MANFRED described protection as by far the most important ecosystem service of mountain forests, and identified eco-engineering as the response where climate change places that function in jeopardy through increased vulnerability to rockfall, avalanche and landslide \citep{jandl_management_2013}.

That objective can be priced. A replacement-cost framework harmonised for the Alpine Space values the protective service by the engineered structures that would otherwise be required, designed to European standards and net of the management expenditure needed to sustain protective capacity. Applied, in a single demonstration case, to a ten-hectare Scots pine stand in the Italian Alps, it returned a stand value of roughly €304{,}000, or €30{,}440 per hectare, equivalent to €950 per hectare per year over a twenty-five-year horizon. The forest was found sufficient to meet local safety requirements without supplementary artificial measures \citep{accastello_asforesee_2019}. The corresponding avoided-risk calculation for a Valais protection forest places annual rockfall risk five times higher without the forest than with it \citep{moos_climate_2021}. Whichever of the two accounting conventions is preferred, the protective value dominates the timber value by a wide margin, and it is that value which fire places at risk.

A management problem follows, however, because protective capacity is not a fixed attribute of a stand but a function of structure that must be actively maintained, and the structures that maintain it are elevation-contingent under climate change. Simulation along a South Tyrolean elevation gradient projects severe declines in avalanche protection in montane forests under strong warming, with drought-driven biomass losses reaching 97\% in dry inner valleys, while subalpine stands hold or improve their rockfall protection \citep{hillebrand_contrasting_2023}. Silvicultural effort therefore needs to be concentrated where the protective function is most exposed, which is not where the forest is most productive.

Multifunctional objectives also conflict in ways that fire sharpens. Scenario modelling of Austrian forests across the century finds stem biomass declining between 9\% and 29\% under harvest-oriented strategies and rising by 34\% under a reduced-harvest, conservation-oriented strategy, and long-lived wood products emerge as the option best balancing sequestration, economic return and disturbance risk \citep{jandl_strategies_2018}. The reduced-harvest option maximises stored carbon and simultaneously maximises standing fuel, and increasing forest cover extends the area over which fire is possible \citep{gerberding_mapping_2025}. Carbon policy and fire policy pull in opposite directions in Alpine forests, and no Alpine study reviewed here evaluates that trade-off explicitly.

If protective structure is the managed objective and fire is an emerging threat to it, then species choice becomes a fire management decision, and it is one Alpine forestry is already making for other reasons.

The direction of travel appears to be set by climate rather than by fire. Simulation of Swiss Alpine landscapes places climate-driven compositional shift ahead of management as the determinant of future biomass and species distribution \citep{schumacher_relative_2006}; participatory modelling of altitudinal vegetation belts in Switzerland anticipates upward shifts requiring site-type reassessment; and observation confirms movement already underway, with treeline advance averaging 10\,m per decade and four fifths of new establishment infilling below the former limit \citep{frei_evidence_2023}. Species suitable for a site today are not necessarily those suitable in fifty years, and Alpine silviculture is accordingly moving towards mixtures and towards provenances drawn from warmer analogues.

Fire cuts across that trajectory in a specific way. The species favoured by warming and drought at low and mid elevations, and those that colonise after fire, are largely not the species that currently deliver protective structure. Post-fire recruitment in the Rhone valley favoured birch, aspen and willow over the pre-fire larch, spruce and Scots pine \citep{moser_potential_2010}; a Valais protection forest is projected to shift from beech- to pubescent-oak-dominated under warming, with rockfall risk rising as it does \citep{moos_climate_2021}. Pioneer broadleaves establish faster than the conifers they replace but hold rock less effectively at the stem diameters that matter, and the transition period is precisely the consequence window described above.

Where fire resistance can be built into species choice, the evidence is reasonably firm and consistent across methods, since stand-level mortality simulation and stem-disc heating converge on one ranking (Section~\ref{sec:forest}) \citep{dupire_vulnerability_2019, mutterer_bark_2026}. Larch and Scots pine, two of the three species that ranking places highest, are already major components of Alpine protection forest, which is a more favourable starting position than the fire-intolerance of the range as a whole would suggest.

The MANFRED programme also evaluated for the Alpine Space what sits alongside species choice. Provenance is a lever in its own right: long-term trials for Norway spruce indicate that potential productivity can be substantially raised, and risk over the stand lifetime reduced, by matching provenance to anticipated future site conditions rather than to present ones \citep{jandl_management_2013}. Browsing, by contrast, is a constraint, and the programme was blunt about it: adaptive silvicultural concepts can only be implemented where browsing damage from high ungulate densities is held to acceptable levels, and at the densities then prevailing, deer were inhibiting the establishment of species-rich forest \citep{jandl_management_2013}. Species diversification is the central recommendation of Alpine adaptive management and ungulate management is its precondition, which places a wildlife-management decision upstream of a fire-resilience outcome.

The unresolved tension is, therefore, with the other objectives silviculture is asked to serve simultaneously. Scenario analysis for Austria finds reduced-harvest strategies maximising carbon storage while building standing biomass \citep{jandl_strategies_2018}; regeneration-oriented framings of mountain forestscape management emphasise moving beyond the maintenance of existing conditions towards actively restoring degraded systems \citep{majlingova_sustainability_2025}; and the harvesting operations through which any of this is implemented are technically and economically constrained in steep terrain \citep{dinca_timber_2025}. No Alpine study reviewed here optimises species selection jointly for protective function, carbon, timber and fire resistance, and until one does, the trade-off is being made implicitly at every planting decision.

Alpine fuel load can be reduced mechanically, by grazing or by fire, and the evidence favours the first two in this setting.

Mechanical treatment has a justification independent of fire. Thinning and shrub removal are recommended in the drought literature because reducing stand competition for water improves carbon acquisition and lowers mortality risk in the remaining stems \citep{giuggiola_droughtinduced_2010}; the same operations reduce surface and ladder fuel and interrupt canopy continuity. In inner Alpine valleys, where drought-driven Scots pine mortality is the standing forest health problem and where fuel accumulation is the emerging fire problem, these are close to the same intervention, and the case for it does not depend on accepting a projection of future fire.

Grazing reaches further, because it carries the broadest policy leverage. Systematic review of the herbivore literature concludes that fuel reduction by grazing animals is a promising route to mitigating wildfire where land abandonment has driven shrub and tree encroachment, with effectiveness contingent on herbivore type, density and feeding behaviour. The same review finds that agricultural and forest policy could support it far more effectively than they do, both by regulating and promoting herbivory in forest areas from which animals have traditionally been excluded to prevent damage, and by directing agri-environmental subsidy towards extensive grazing and pastoralism \citep{rouetleduc_effects_2021}. Landscape simulation quantifies the gain from embedding fuel treatment in land-use policy rather than operating it alone: prescribed burning by itself reduced subsequent wildfire by up to 36\% over thirty years, but the leverage of treatment reached unity (one hectare of wildfire avoided per hectare treated) only in combination with high nature value farmland restoration and fire-smart forest conversion \citep{pais_optimizing_2023}.

The effect sizes available for grazing are large, and all of them come from outside the Alps. Managed grazing has reduced rate of spread by more than 60\% in pure grass communities and by over 50\% in mixed grass and shrub, and Mediterranean studies report burnt-area reductions between 25.9\% and 60.9\% \citep{xu_how_2026}. The mechanism is patchiness rather than uniform fuel removal, since grazing interrupts horizontal continuity, which makes the instrument behave more like a distributed set of small breaks than like a thinning. Its response to intensity is non-linear and reverses: removal of less than about half the photosynthetic leaf area reduces fuel and favours fire-resistant species, whereas heavier use, particularly during drought, exacerbates erosion, admits flammable annuals and can leave a continuous homogeneous fuel bed more flammable than the one it replaced \citep{xu_how_2026}. On erodible Alpine slopes above assets that reversal has a second cost, because the same overgrazing that raises flammability also removes the ground cover that arrests small blocks \citep{asteriou_assessing_2025}. The threshold separating the two regimes has not been located for any Alpine pasture, and the seasonal question compounds it, since effectiveness in high-altitude systems is greatest in the drier spring and autumn windows rather than in the humid summer when Alpine stock is at its highest elevation \citep{xu_how_2026}. The only Alpine evidence is historical and correlational. In the Valais, livestock density in forests was associated with fewer fires while wood pasturing remained widespread before the late 1950s, whereas after 1956 high livestock densities coincided with more fires. Those densities lay in the driest eastern part of the canton, where dryness rather than fuel removal is the likelier explanation \citep{zumbrunnen_human_2012}.

Where treatment sits in the landscape matters as much as whether it happens. Landscape modelling of a transboundary Iberian biosphere reserve over 2010--2050 found that restoring high nature value farmland delivered a present-value suppression-cost saving of roughly €3.9 million, and that combining it with fire-smart forestry cut total societal cost (suppression plus ecosystem service damage) from €25.7 million under business as usual to €7.1 million. Fire-smart forestry applied alone, by contrast, \emph{raised} suppression costs; the saving came from the agricultural-landscape component \citep{lecinadiaz_incorporating_2023}. For the Alps, where abandonment of agricultural land is the proximate cause of fuel accumulation, this is the most directly relevant economic result available, and it argues that the effective lever is agricultural policy rather than forest policy alone.

Empirical support for that argument now exists at national scale, if not for the Alps alone. Across about 49{,}000 Italian fires larger than a hectare between 2007 and 2017, analysed in 602 agro-forest districts spanning Alpine, temperate and Mediterranean climate regions, districts with more active land governance showed smaller burnt shares of ecosystem services and of the wildland--urban interface and lower forest burn severity. The association held once climate, fire weather and simulated flammability had been accounted for, while land abandonment raised all three impacts \citep{spadoni_active_2023}. Governance was proxied by sustainable forest management and certification, extensive agro-pastoral activity, protected areas, and rural development and LIFE funding. Its effect was moderate beside climate and flammability, and not uniformly monotonic: severity fell with the governance of rural and marginal areas up to the median value and rose beyond it, a reversal the authors relate to intensive harvesting.

The effectiveness of treatment is also conditional on the fire weather under which it is tested, and the conditionality runs against the direction the Alps are moving. A randomised, replicated experiment that subsequently burned in a wildfire found combined thinning and prescribed burning reducing the probability of tree mortality by 0.60, crown volume consumed by 86\% and bole char height by more than seven metres relative to untreated controls; but the mortality benefit was maximised at intermediate values of the burning index and declined under the most extreme fire weather, even as crown-fire metrics remained responsive \citep{brodie_forest_2024}. Fuel treatment buys most where danger is moderate, and land governance shows the same dependence at landscape scale, with its mitigating effect strongest in warm, dry and flammable districts and weakening as the fire weather during the events became more severe \citep{spadoni_active_2023}. Since Alpine danger is currently moderate and rising \citep{gerberding_mapping_2025, miller_climate_2024}, the window in which treatment is most cost-effective is open now and narrowing.

\subsection{Prescribed fire, fuel breaks and post-fire restoration}
\label{sec:fuelbreaks}

For prescribed fire, fuel breaks and post-fire restoration alike, the Alpine evidence is thinner than the confidence with which each instrument is discussed, and prescribed-burning practice is unevenly distributed. Prescribed burning is applied at scale in France, where 1{,}000 to 3{,}000\,ha are treated annually in the Alpes-Maritimes by a trained technician community with shared protocols, but in the Mediterranean south-east rather than in the Alpine region proper, and for grazing recovery as much as for hazard reduction \citep{muller_forest_2020}. The technique, the training structure and the institutional model therefore exist within an Alpine country without having been transferred to Alpine conditions.

Prescribed burning has been treated as a poor fit for the Alps, and the structural reasons are real. The dominant species carry little fire adaptation, with silver fir, Norway spruce and especially beech ranking lowest in experimentally measured bark insulation \citep{mutterer_bark_2026, dupire_vulnerability_2019}. The burn window is narrow and coincides with the spring anthropogenic fire peak \citep{muller_forest_2020}, and escape consequences are severe because the slopes that need treating are those above assets. In Austria the obstacle is also legal, since lighting fires in forest is prohibited under the Forestry Act and prescribed burning is not considered a management option \citep{neumann_fuel_2022}. Even the assumption that low-intensity burning is hydrologically benign is undermined by the finding that the strongest post-fire water repellency occurred in low-severity plots, driven by abundant uncharred organic matter rather than by intense heating \citep{negri_wildfire_2026}.

One Alpine trial provides encouraging evidence for prescribed burning in Scots pine. Eighteen plots at a south-facing site at 1600\,m showed roughly halved fuel loads over one year, with little measured change in plant hydraulics or soil properties \citep{gamba_prescribed_2025}. The result supports further trials in comparable stands, while leaving operational-scale safety, other species and longer-term responses unresolved.

Evidence from the soil biota points the same way, though from outside Europe. In paired southern Appalachian watersheds, the bacterial and fungal communities of repeatedly prescribed-burnt stands were statistically indistinguishable from those of fire-excluded stands, while wildfire-affected and partially harvested stands differed significantly from both; ectomycorrhizal abundance held at 63 to 66\% under prescribed fire and exclusion but fell to 38\% after wildfire, with saprotrophs rising correspondingly \citep{snyder_taxonomic_2026}. The comparison is confounded by differing times since disturbance and by elevation, and the forest type is not Alpine, so it cannot be carried across as a result. What it does is identify the soil microbial community as a measurable endpoint on which prescribed fire and wildfire separate cleanly, which is the kind of endpoint the Alpine prescribed-burning question currently lacks.

A European mountain range has meanwhile moved past the trial stage, and its setting resembles the Alps more closely than the Mediterranean cases usually cited. The Aran valley in the Spanish Pyrenees, some 634\,km$^2$ at a mean altitude of 1{,}840\,m and a mean slope of 27\degree{}, under an Atlantic rather than a Mediterranean climate and with a winter-dominant fire season running from January to March, adopted a strategic fire management plan unanimously in 2022. The plan is organised around an explicit target for how much fire the landscape should carry, set at 0.66\% of natural land cover annually against a historical fire rotation of about 350 years, and its first implementation burnt 327\,ha over fourteen working days inside a pre-planned 346\,ha management area \citep{oliveres_path_2025}. Its winter and early spring burning windows, its transhumance history in which shepherds used fire as a routine tool, and its deep organic soils, in which fire can smoulder and reignite, all transfer to the Alps directly, the last because it is the suppression problem that Alpine lightning fires present.

The Pyrenean case also shows what adoption requires, which is not primarily ecological evidence. The plan rests on a community participation process with municipal representatives, on eight formally agreed goals spanning severity reduction, emergency resilience, extensive livestock farming and biodiversity, and on landscape fire succession modelling to locate the treatment. Its authors are candid that economic sustainability has not yet been assessed, that managed fire can escape, and that explicit wildfire use under technical and strategic direction remains far from common European practice \citep{oliveres_path_2025}.

The Alpine trial and the Pyrenean programme address different questions: ecological response at one experimental site and implementation across a managed landscape. Together they justify evaluation that combines species response, soil monitoring, escape risk and local participation. Neither establishes a generally suitable burning prescription for the Alps.

Outside mountains, the standard landscape instrument for limiting fire size is linear, and its terminology distinguishes forms by how much vegetation survives the treatment. A firebreak is a strip cleared to mineral soil; a fuelbreak retains herbaceous and shrub cover while reducing load and vertical continuity; a shaded fuelbreak retains the tree canopy and treats only the understorey and ladder fuels \citep{ascoli_firebreak_2018}. The distinction is not pedantic, because it determines whether the instrument is compatible with an Alpine protection forest.

On that criterion the mountain form already exists and is recommended. Linear fuelbreaks and shaded fuelbreaks are held to be better suited than cleared firebreaks precisely to mountainous territory with steep slopes and rough morphology, because retaining vegetation cover limits soil erosion; and where terrain is very steep, ridgelines and valley bottoms are the preferred siting \citep{ascoli_firebreak_2018}. Recommended widths run from 30 to 100\,m for primary networks and from 3\,m to over 50\,m for secondary and tertiary ones, with siting criteria that already include maintenance cost, accessibility along roads and firefighting considerations. European Mediterranean practice adds detailed prescriptions at the interface, differentiating an inner zone in which trees are removed from an outer zone in which they are retained but pruned and thinned, with shrub phytovolume thresholds and crown-separation distances \citep{rigolot_recommendations_2003}. The Alpine literature reviewed here contains no equivalent design guidance calibrated to Alpine fuels, slopes and protective requirements, which is a gap in application rather than in principle.

How much such a network could deliver in the Alps is nonetheless limited, and the limits can be quantified. Linear barriers are, to begin with, less effective than their prominence in policy suggests. A recent synthesis finds firebreaks stopping only 28\% of large fires in southern California between 2017 and 2020, some 73\% of Western Australian wildfires breaching them fully or partially, and a proposed national network in Portugal projected to reduce fires above 500\,ha by 13\% while still being breached in more than a third of such events, with firebrand transport across the gap the dominant failure mechanism \citep{zong_effectiveness_2026}.

Where they do work, the reason inverts the usual understanding of what these features do, and it points to the constraint that binds hardest in mountains. Analysis of 563 fire--fuelbreak intersections from 220 wildfires in Andalusia found containment in 46.9\% of cases overall, but with an extreme dependence on whether suppression was operating on the line: 6.66\% with no suppression, 15.9\% with aerial resources alone, and 76.74\% with ground crews working it \citep{ortega_modeling_2024}. A fuelbreak is therefore a platform for ground suppression rather than a barrier in its own right. Its value is contingent on crews reaching and holding it, and in Alpine terrain the same slope, access and exposure conditions that motivate building one are those that prevent ground crews from occupying it. The terrain-conditioned suppression-difficulty framework discussed in Section~\ref{sec:governance} is the natural instrument for identifying where that contingency binds, and it has not been applied to Alpine fuelbreak siting.

Geometry compounds this. Effectiveness in the Andalusian data depended strongly on the angle at which fire met the line, with near-parallel intersections performing best and perpendicular intersections succeeding in only 5.6\% of cases \citep{ortega_modeling_2024}. Alpine fires run upslope along the paths whose length predicts large-fire likelihood \citep{conedera_linking_2024}, so a contour-aligned break meets that run perpendicularly. Ridgeline siting \citep{ascoli_firebreak_2018} partly resolves the tension, since a ridge intercepts the run at its terminus rather than across its flank, but it also places the treatment at the position most exposed to wind and least accessible to ground crews.

Even a well-sited network would then meet the obstacles of ownership and finance. Large and long linear networks are held to be poorly suited where the wildland--urban interface is widespread and where forest property is largely private with small cadastral parcels \citep{ascoli_firebreak_2018}, a description that fits much of the Alpine arc. The Catalan experience quantifies the resulting failure mode: prevention subsidies reached under 4\% of designated strategic area and left an average minimum distance of 287\,m between funded parcels, so the spatial continuity the scheme required was never achieved \citep{gorrizmifsud_current_2026}. Nor does the financing model transfer. Optimisation of a 3{,}538\,km network in the north-western United States found around 97\% of segments requiring thinning with pile burning, underwritten by harvest revenue of the order of US\$63--75 million \citep{ager_optimizing_2023}; Alpine harvesting in steep terrain is technically demanding and economically marginal \citep{dinca_timber_2025}, so that revenue is largely unavailable.

A residual trade-off remains specific to protection forests, and using the shaded form does not eliminate it. Even a treatment that retains canopy reduces stem density and removes the lying material that contributes to barrier effect, and both govern rockfall arrest \citep{costa_biological_2021, moos_climate_2021}; against a replacement-cost valuation of Alpine protective capacity of the order of €30{,}000 per hectare \citep{accastello_asforesee_2019}, the acceptable degree of thinning is an empirical question with a monetised answer that nobody has computed. Establishing how much understorey and stem removal a protection forest can absorb before its protective function degrades faster than its fire hazard falls is, on the evidence assembled here, the single most tractable unanswered question in Alpine fuel management.

One line of support for these instruments comes from the vulnerability side rather than the suppression side, and it applies at a different scale. The expert weighting described in Section~\ref{sec:risk} placed protective strips and firebreaks at the head of the building-environment indicators, and the judgement being expressed there concerns the building rather than the landscape. A break adjacent to a structure functions as a defensible margin, which is not the containment function the effectiveness literature evaluates, and the terrain that makes a landscape-scale line impossible to hold does not prevent a cleared margin around a house. Separating the two uses resolves an apparent contradiction between an instrument that performs poorly at landscape scale and one that practitioners rank first at the scale of the exposed asset.

The broader implication, in our reading, is that Alpine fuel management is likely to be predominantly area-based rather than geometric, with linear elements confined to ridgelines, existing road corridors and the interface. The instruments the evidence supports are grazing that reduces load across whole abandoned slopes \citep{rouetleduc_effects_2021}, thinning justified simultaneously by drought stress and by fuel \citep{giuggiola_droughtinduced_2010}, prescribed burning where the stand tolerates it \citep{gamba_prescribed_2025}, and prioritisation targeted at the terrain positions from which large fires actually develop \citep{conedera_linking_2024, krsnik_spatial_2024}.

Post-fire management can influence recovery, but receives less research attention than prevention and active fire. A bibliometric analysis counted approximately 3{,}976 post-fire publications, compared with 5{,}802 on prevention and 7{,}041 on active fire \citep{goncalves_fostering_2025}. These counts indicate uneven attention, although publication volume alone does not measure the practical adequacy of the evidence.

This deficit is also institutional, and it has now been characterised at European level. Post-fire management is the least developed component of the fire management cycle across the European Union, and the obstacle identified is an implementation gap between available science and operational practice rather than an absence of tools, since stabilisation measures, restoration techniques and monitoring protocols all exist; what is missing is common assessment procedure, multidisciplinary operational teams, long-term funding and knowledge exchange between member states \citep{vieira_postfire_2026}. That diagnosis was drawn from Greek, Portuguese and Spanish experience, which makes its transfer to the Alps a question rather than a conclusion, but the Alpine case is the harder one, because the consequence the intervention is meant to forestall is gravitational rather than agronomic and the window in which it must act is set by the recovery of protective structure (Section~\ref{sec:forest}).

The Alpine evidence that exists is critical of current practice. Analysis of short-term recovery after a large mixed-severity fire in the Western Alps concluded that the restoration techniques most often applied are poorly matched to the ecological context and to the consequences of global change, and that intervention can either facilitate recovery or trigger further degradation \citep{mantero_short-term_2023}. The biological constraint is drought: post-fire recruitment at low and mid elevations in the Central Alps requires precipitation well above the long-term average, and it is moisture rather than fire intensity that governs seedling density \citep{moser_potential_2010}. At the upper margin the constraint may be absolute, with no meaningful mountain pine re-establishment recorded over eighty years on burnt subalpine slopes \citep{degiuli_arrested_2024}.

Post-fire salvage decisions should account for the protection supplied by surviving and dead stems. Evidence from burnt southern Alpine beech stands favours retaining suitable deadwood, while a windthrow study demonstrates the barrier effect that lying logs can provide \citep{maringer_temporal_2016, costa_biological_2021}. Longer-term Tatra observations also show trade-offs between deadwood retention and regeneration diversity \citep{seben_impact_2024}. These findings inform management, but do not replace site-specific assessment of unstable trees, future fuels and downslope protection.

\subsection{What differs in mountains}

Where forests protect downslope assets, fuel management must retain enough structure to intercept rockfall and stabilise slopes. Alpine studies value this service and show its dependence on stand condition \citep{accastello_asforesee_2019, moos_climate_2021}. Treatment can alter that protection, but the trade-off depends on which stems, understorey and deadwood are removed. It is not an inevitable consequence of every reduction in fuel.

The history of those fuels also matters. Abandonment of grazing and extraction has contributed to forest expansion and continuity in the Alps, while evidence from the Apennines shows that present composition reflects a much longer interaction between land use, climate and fire \citep{bebi_changes_2017, garbarino_contrasting_2020, moralesmolino_longterm_2021}. These findings support coordination between forestry and agricultural policy. They do not exclude a role for suppression or imply that abandonment is specific to mountains.

Management options depend in turn on species and site. Fire-resistance experiments distinguish relatively vulnerable beech, fir and spruce from better-insulated larch and Scots pine \citep{mutterer_bark_2026}. A single Alpine prescribed-burning trial in Scots pine provides encouraging evidence, but does not establish suitability across forest types \citep{gamba_prescribed_2025}. The appropriate treatment therefore depends on both expected fire response and the protection required below the stand.

After fire, delayed mortality and regeneration determine how long that protection is reduced. Southern Alpine beech studies show losses lasting decades, while one subalpine massif showed little tree re-establishment over eighty years \citep{maringer_temporal_2016, gehring_shallow_2019, degiuli_arrested_2024}. Dead stems may retain some protective function before decaying or falling. Their condition must therefore be assessed directly when deciding whether post-fire salvage would improve or reduce protection.

% =======================================================
\section{Risk, Exposure and Vulnerability}
\label{sec:risk}

\subsection{From hazard to risk: the missing terms}

A single-hazard wildfire assessment asks where fire is likely and how large it might be. In the Alps that is the wrong terminal question. Burnt area is small and weakly related to consequence; the exposed asset usually lies outside the burn perimeter but beneath it; and susceptibility, correctly defined, is a static probability of experiencing fire given the intrinsic characteristics of the territory \citep{trucchia_defining_2022}, a legitimate quantity, but separated from risk by an unusually wide gap in mountains.

Risk combines hazard, exposure and vulnerability. Alpine research has made progress on all three, including a building-vulnerability index, but the reviewed studies do not provide an integrated assessment across the whole arc \citep{gerberding_mapping_2025, papathomakohle_towards_2025}. A hazard surface needs additional information about exposed elements and their potential losses before it can be interpreted as risk.

Both missing terms have been developed in a separate literature with its own instruments and data requirements, namely indicator sets weighted by expert judgement where damage records are absent \citep{papathomakohle_towards_2025}, indices calibrated on surveyed damage where such records exist \citep{papathomakohle_wildfire_2022}, and, for the ecological component, continental mapping that integrates lost ecological value over the recovery period \citep{arrogantefunes_mapping_2024}. Neither term has been joined to an Alpine hazard surface. This is the gap that most directly separates the present Alpine literature from a usable risk assessment, and closing it is a matter of combining layers that already exist rather than of improving any one of them (Section~\ref{sec:governance}).

\subsection{Exposure: habitats, assets, people and seasonality}

Exposure in mountains is neither static nor confined to the burn perimeter, and its ecological component is the one already measurable at continental scale, where the measurement is far from marginal. Of the 1{,}079{,}538\,ha burnt in the European Union in 2025, the largest total in the European record, 424{,}023\,ha lay inside Natura 2000 sites: 39\% of the total, the highest share ever recorded, three times the 2024 figure and twice that of 2023, with Spain, Portugal and Romania accounting for two thirds of it \citep{jrc_advance_2026}. European burning does not incidentally touch the protected-area network; it falls inside it in roughly two fires out of five.

For Austria, the 2025 report maps four fires totalling 434\,ha, including 420.6\,ha in Natura 2000 sites and 334\,ha classified as agricultural land \citep{jrc_advance_2026}. These national figures do not establish where the fires lie within the Alpine perimeter, nor whether the agricultural land was abandoned. They illustrate the need to check geography and land-cover categories before treating national totals as Alpine exposure estimates.

We therefore argue that exposure analysis must follow the cascade rather than the fire. The literature that does this best is the vulnerability rather than the hazard literature: a participative multi-criteria approach in south-eastern France, distinguishing exposure, sensitivity and adaptive capacity across populations, ecosystems and infrastructure, identified secondary vulnerability hotspots and locations not yet affected by hazard where prevention was nonetheless warranted \citep{riviere_mapping_2023}. Exposure is also dynamic. Human presence in Alpine forests is tourism-driven and seasonal \citep{pede_mountain_2022, abrham_wildfires_2024}, and the leading European ignition model identifies distance to the wildland--urban interface as the strongest human predictor while conceding that its human covariates are static and should be made to vary with commuting and seasonal recreation \citep{gelabert_assessing_2025}. Transboundary simulation in the Saxon--Bohemian borderland illustrates the same point institutionally, with modelled burn probability and flame length at settlements differing across a national border that the hazard does not recognise \citep{avouris_exposure_2025}.

Compound structure enters the exposure picture at more than one point, and conflating those points has caused confusion. Preconditions compound, since European fire activity typically follows moderate but coincident hydro-meteorological and land-surface anomalies rather than single extremes \citep{miller_compounding_2025}. Disturbances compound, since drought, windthrow and bark beetle outbreak load the fuel bed before any fire arrives, a pathway explicitly identified but not represented in Alpine hazard projections \citep{gerberding_mapping_2025, bebi_changes_2017}. Consequences compound as well, in the cascade described above, for which a research roadmap now exists that treats post-fire and frozen-soil dynamics together and calls for multi-institutional monitoring networks and for governance linking geoscientists, engineers and policymakers \citep{baselt_compound_2026}. The exposure of linear infrastructure to that consequential form has been quantified, although not in Europe. Across California, the likelihood that the historic mean coincidence of intense rainfall over burnt ground would be exceeded rises from 38\% to 68\% and 73\% under moderate and high emissions pathways, and for the upper decile of events from 34\% to 90\% and 91\%. Natural gas transmission is thus projected to face a compound exposure substantially greater than the sum of its exposure to either driver \citep{moftakhari_increasing_2019}. The method is transferable even though the coefficients are not, and the Alpine equivalent, roads, rail and pipelines beneath burnt slopes in confined valleys, has not been attempted.

\subsection{Vulnerability of the built environment}

Vulnerability has been constructed for the Alps from expert judgement, and the exercise shows what is missing. An Austrian index scores buildings in the wildland--urban interface on twenty-one indicators grouped as building environment, building exterior and building characteristics, weighted by pairwise comparison across a national and an international expert panel and applied to twenty-three buildings in a Carinthian settlement on the southern slope of the Dobratsch \citep{papathomakohle_towards_2025}. Both panels ranked the presence of protective strips or firebreaks first among the six environment indicators (weights of 0.286 and 0.198), the distance of tree crowns from the building first among the seven exterior indicators (0.258 and 0.212), and roof material first among the eight building indicators (0.203 and 0.164). The weights rest on expert judgement rather than on damage records because no Austrian wildfire damage record exists from which to derive them.

A Greek index calibrated on 423 buildings after the Attica fire ranked roof material, structural type and slope highly, while the Austrian expert panels assigned relatively little weight to slope \citep{papathomakohle_wildfire_2022, papathomakohle_towards_2025}. Roof characteristics have a clear relationship to ember and heat exposure \citep{naser_vulnerability_2025}. However, agreement on an indicator does not establish that its numerical weight transfers between settings. The difference in slope weighting also reflects contrasting evidence and methods, not a controlled test of regional transfer.

Where fire does reach structures directly, the wind interaction that terrain amplifies is also what drives damage: review of fire--wind simulation identifies the fanning effect of wind, the combination of wind direction with slope angle, and the resulting flow around buildings as the processes governing wildland--urban interface losses, and notes that disturbance interactions affecting structures have been comparatively neglected \citep{ghodrat_existing_2021}. Alpine settlement sits in valleys where slope and channelled wind combine by default.

\subsection{What differs in mountains}

Risk assessment in mountains must consider direct burning alongside hazards that travel beyond the perimeter. An Alpine debris flow and rockfall observations show how fire can alter conditions above exposed settlements and infrastructure \citep{laporta_postwildfire_2025, melzner_brief_2019}. Such consequences also occur outside the Alps; their relevance here follows from the location of assets below steep, forested slopes. Exposure mapping must include the affected slopes, channels and runout zones.

The people using those assets also change through the year. Alpine tourism studies describe shifts in seasonal demand under warming, suggesting a need to represent visitors as well as resident populations \citep{pede_mountain_2022}. This requirement is not unique to mountains, but becomes consequential where seasonal access, evacuation and fire weather coincide.

Once exposure is located, vulnerability still requires local evaluation. Austrian expert-based and Greek damage-based building indices assign different importance to slope, despite both identifying roof material as relevant \citep{papathomakohle_towards_2025, papathomakohle_wildfire_2022}. Agreement identifies a useful candidate indicator; it does not establish transferable weights. Protective and ecological losses also require separate measures, because stand recovery, habitat change and direct building damage are different outcomes \citep{gehring_shallow_2019, rey_effects_2019, arrogantefunes_mapping_2024}.

% =======================================================
\section{Governance, Policy and Economics}
\label{sec:governance}

\subsection{Suppression, preparedness and the practitioner constraint}

Alpine suppression operates at once under terrain-driven behavioural unpredictability \citep{sharples_overview_2009, melzner_brief_2019}, low accessibility in steep and remote ground \citep{muller_forest_2020}, and dependence on shared aerial resources whose availability is correlated across the arc precisely when synoptic-scale drought and vapour pressure deficit elevate danger everywhere at once \citep{miller_compounding_2025}.

What those constraints mean in practice is a suppression system organised around water logistics and vertical distance rather than around a frontal assault. Where vehicles can reach, ground attack works from tank tenders in combination with ponds, hydrants and long hose lines, with independent pumps inserted to overcome large differences in height; where they cannot, initial attack on ground fire is made with rakes and beaters, and subsurface hot spots are reached with backpack pumps and high-pressure lances. Vehicle specification follows the road rather than the fire, since Italian forest-fire units operate all-wheel-drive pickups sized for narrow Alpine forest roads rather than the heavier tankers used on open ground, and sprinkler lines wetting vegetation along fuel breaks are used in the southern Alps, particularly at night. Tactical or counter-firing during an incident, which is distinct from prescribed burning as a prevention measure, is practised mainly in France and Italy and is uncommon elsewhere in the range \citep{muller_forest_2020}.

Aerial support is a structural dependency in this system rather than an adjunct, and its function differs from the Mediterranean one. Instead of delivering retardant along a long front, helicopters in Austria and Bavaria increasingly ferry water to mobile basins placed near the fire from which ground crews lay hose, and every Alpine country maintains units trained specifically for bucket work and winch insertion \citep{muller_forest_2020}. The constraint being addressed is vertical rather than horizontal, and this corrects an assumption that is easy to import from Mediterranean practice. Forest road density in the Alpine region is high by comparison with Mediterranean or Scandinavian Europe \citep{muller_forest_2020}, so the accessibility problem is not an absence of roads but the slope, exposure and vertical separation between the road and the fire, which is the quantity the terrain difficulty index formalises \citep{cardil_assessing_2025}.

Hazards to which Alpine crews are exposed are correspondingly different. Rockfall and fall risk are ordinary features of the working environment, burning material rolling downslope ignites new fires below crews working above them, and in sectors of the range where the two world wars were fought, buried munition remains a recognised hazard; where terrain is steep enough, mountain rescue teams are integrated into the action force for the purpose of securing the firefighters themselves \citep{muller_forest_2020}. Confined topography adds a behavioural hazard on top of these, since canyons can produce sharp increases in fire intensity without any accompanying increase in wind \citep{muller_forest_2020, hayajneh_wind_2026}. A suppression doctrine written for terrain in which the principal danger is the fire front does not describe this working environment.

Organisationally the Alpine system also differs, and this is the part least represented in the modelling literature. Dense networks of volunteer brigades in Austria and Switzerland, together with a population that reliably reports smoke by telephone, mean that most Alpine fires are attacked within a short time of ignition \citep{muller_forest_2020}. We take this to be part of the explanation for a fire-size distribution dominated by events under a hectare, although the inference has not been tested: the same fire weather in a range without that response density would plausibly produce a different distribution, and no Alpine hazard model contains a term for it. It is also a capability with a demographic exposure of its own, since it rests on voluntary participation in valley communities that are ageing and contracting (Section~\ref{sec:economics}).

A quantitative framework for the accessibility constraint now exists, though it has not yet been applied in the Alps. Combining a terrain difficulty index built from accessibility, penetrability, line-construction difficulty, mean slope and road proximity with a fire behaviour index yields an initial attack assessment whose highest class carries a 90\% reduction in the odds of successful initial attack relative to the lowest, with mean burnt area rising from 18.5 to 1731\,ha across that range, and with response time interacting with terrain difficulty to further reduce containment probability \citep{cardil_assessing_2025}. Since the framework has been validated outside its region of origin, the Alpine gap is one of application rather than of method, which converts an apparently intractable evidence gap into a tractable one.

Where suppression cannot be fast everywhere it must be fast in the right places, and the right places are identifiable in advance from terrain and fuel connectivity, which is the operational argument for targeted placement of pre-suppression infrastructure such as water points for aerial firefighting \citep{conedera_linking_2024}.

The interaction between suppression capability and fuel treatment deserves particular attention in a range whose danger distribution is shifting. Fuel treatment reduces the severity a suppression force has to confront, but its benefit for tree mortality is greatest at intermediate rather than extreme fire weather and declines under the most severe conditions \citep{brodie_forest_2024}; suppression success, conversely, collapses with terrain difficulty largely irrespective of fuel \citep{cardil_assessing_2025}. In accessible Alpine terrain under moderate danger the two are substitutes, and the cheaper can be chosen. On the steep, remote slopes where fuel connectivity and upslope path length are greatest \citep{conedera_linking_2024}, suppression is weak and treatment is the only lever that acts before the fire; but these are also the sites where mechanised treatment is most costly \citep{dinca_timber_2025}, which is why the pastoral and silvicultural instruments of Section~\ref{sec:forest} matter disproportionately there.

Alpine fire management currently spends several times more on response than an integrated prevention framework would cost, a ratio quantified in Section~\ref{sec:economics}, and econometric analysis of Italian regions reaches the same qualitative conclusion, that upstream prevention is the more efficient allocation under projected climate intensification \citep{michetti_forest_2019}. What no Alpine study provides is the quantity that would settle the allocation: the marginal effectiveness of a euro spent on suppression capability against a euro spent on fuel treatment, conditioned on terrain.

The strategic question is how much of the projected increase management can absorb, and the continent-scale projections introduced in Section~\ref{sec:intro} bound the answer: sustained investment in capacity absorbs most of the climate-driven increase in burnt area, but not all of it, and least of all in the foothill and medium-elevation terrain where the increase concentrates \citep{billing_future_2026}. For the Alps, investment in capacity is therefore the largest single lever available within the region's control rather than a marginal adjustment, yet it will not be sufficient on its own in the terrain where the increase is projected to concentrate.

\subsection{Prevention: placement, early warning and integrated fire management}

The instruments of Alpine prevention were set out in Section~\ref{sec:forest}; the governance question is whether they are deployed where they are needed. Fire-smart management is already the operational focus in Mediterranean regions, and the argument for building a pre-fire capability in the currently less affected Alps rests on anticipation rather than on experience \citep{gerberding_mapping_2025}, which makes the allocation problem political as much as technical.

The Alpine evidence on allocation is, however, discouraging. Of 97 nature-based solutions for climate adaptation mapped across the European Alps, most addressed drought or soil instability and aimed to deliver multiple benefits at once, yet they were generally not located within the hotspots of the hazards they claimed to address, nor within hotspots of the ecosystem contributions they aimed to secure, nor where the two coincided \citep{dubo_nature-based_2023}. Measures are being placed where they are institutionally feasible rather than where they are physically needed, and closing that gap requires local climate projections and quantified ecosystem contributions to be integrated into planning, through a macro-regional strategy combined with local stakeholder engagement.

A corrective is available and unusually cheap. Because large-fire likelihood in mountain terrain can be mapped from slope length and fuel connectivity using only a digital elevation model and a forest map \citep{conedera_linking_2024}, and because protective function can be valued per hectare by replacement cost \citep{accastello_asforesee_2019}, the two layers needed to rank candidate treatment sites by expected avoided loss already exist for most of the Alpine arc. What is missing is the institutional mandate to combine them, since the hazard layer sits with fire and forest services and the value layer with the authorities responsible for the assets below.

A transferable architecture for that prioritisation already exists. Multi-criteria analysis of some 238{,}000 forest polygons across Catalonia ranked management zones on fuel hazard, simulated fire behaviour, exposure of settlement and infrastructure, landscape continuity and, most importantly in mountains, accessibility, expressed through terrain aspect, road density and distance to water points and strategic management areas \citep{krsnik_spatial_2024}. The method is directly portable to Alpine conditions, where accessibility is the binding constraint on suppression. Its sensitivity is also instructive: expert weighting designated 20.69\% of the territory as high priority while equal weighting designated 7.98\%, so the participatory step determines the answer as much as the data does.

We take the most tractable near-term gain to be not a new Alpine model but the extension of early warning systems that already exist. The Alps possess dense hydro-meteorological monitoring, mature torrent and avalanche warning services and long institutional experience of gravitational hazard; what those systems lack is a fire term. The concrete proposal in the Alpine literature is to modify catchment runoff parameters automatically for burnt catchments, so that a debris flow early warning system lowers its rainfall thresholds without waiting for a field survey, using thematic layers (catchment and fan delimitation, geology, bedrock outcrop, land use and satellite-derived fire perimeters) that are already available in geographic information systems \citep{tiranti_wildfires_2021}. The Bussoleno event supports the need for exactly this, since the triggering intensity fell below the threshold that an unburnt-basin system would have been watching for \citep{laporta_postwildfire_2025}.

The modelling evidence also bears on how such a system should be designed. Its outputs are most useful when expressed as decision thresholds rather than as probability surfaces, with the asymmetry between missed detections and false alarms treated as part of the specification rather than as a post-hoc consideration; Alpine ignition modelling has already been built this way, deriving dynamic thresholds from combinations of true-positive and false-alarm rates \citep{moreno_1_nodate}, and integrated danger assessment has reported false-alert performance as a first-class result \citep{muller_towards_2020}. In a setting where a missed detection means an unwarned slope above a village, this is not a refinement.

By the same logic, integrated fire management is best specified around indicators that different institutions can each observe and act on. The Alpine post-fire inventory work proposes precisely that: identification of indicators for burn severity and for subsequent geohydrological processes, with systematic documentation procedures, framed explicitly for use within an integrated fire management strategy by policymakers rather than by researchers \citep{melzner_wildfire-induced_2025}. The costed Alpine framework already allocates roughly a fifth of a proposed annual budget to post-fire management and a further tranche to knowledge transfer \citep{muller_forest_2020}, and the compound-hazard roadmap specifies the monitoring networks that would supply the indicators \citep{baselt_compound_2026}. The components have been separately proposed; what is missing is their assembly.

\subsection{Fragmentation, harmonisation and cross-border coordination}

The Alpine wildfire cascade crosses institutional boundaries at every link: prevention sits with forest administrations, suppression with fire services and civil protection, post-fire rockfall and torrent control with geological surveys and infrastructure operators, and the land-use decisions governing fuel continuity with agricultural and planning authorities. The sequence described in Section~\ref{sec:water} passes through all four.

Drivers themselves are nationally differentiated, which is one reason a single Alpine prescription is unlikely to fit. Rural abandonment dominates in France and Italy, while tourist activity dominates in Austria, Germany and Switzerland; negligence is the leading proximate cause throughout \citep{muller_forest_2020}. Prevention that works through agricultural policy in the western and southern Alps must work through visitor management in the north and east.

Beneath the institutional division, the fragmentation reaches the level of basic factual disagreement. Surveying three stakeholder groups across seven Alpine countries found significant differences in perceptions of fire numbers, burnt area and future trend, even where the groups agreed that precipitation, temperature and wind are the dominant natural drivers and that recreation, rural abandonment and owner awareness are the dominant socio-economic ones \citep{hoang_stakeholders_2026}. A hazard whose practitioner community does not share a factual picture is not one that is ready to grow.

How that plays out where citizens have formal standing has been surveyed in one Alpine setting. In the Locarnese region of Ticino, residents combined a high valuation of local amenity and strong support for firefighting services with a diagnosis that the principal driver of risk at the forest and settlement interface is the current level of urban expansion rather than the hazard itself. Despite a planning system built on deliberative instruments, younger, locally rooted and less formally educated respondents reported the least satisfaction with their influence over it, and the availability of popular referenda was judged to foster a disproportionate sense of entitlement over local decisions \citep{bardsley_wildfire_2021}. For a range in which land-use decisions govern fuel continuity, the implication is that public support for suppression capacity is not evidence of public support for the land-use restraint that would reduce exposure, and that formal participation rights do not distribute influence evenly among those who carry the risk.

The same diagnosis emerges from within a single Alpine country when practitioners are asked directly. A structured survey of 79 wildfire experts across seven Italian regions, three of them Alpine, found the institutional structure highly fragmented, with responsibility increasingly distributed among separate organisations, and emergency response persisting as the default strategy despite formal commitments to precaution; local stakeholder participation was reported as displaced by a shift towards extra-local actors, and resources at state level were judged disproportionate to regional capacity. Prevention dominated both halves of the assessment: 46\% of the coded shortcomings concerned prevention deficits, and 58\% of the identified opportunities concerned enhanced prevention \citep{kirschner_governance_2024}. That the same term dominates what is missing and what is possible is the clearest available statement that the constraint is institutional rather than technical.

Fragmentation is also visible one level up, in the European policy framework itself. A review of forty EU sectoral policies across seven clusters found integrated fire management named seven times in the entire corpus, with a substantive mention in a single policy text, and the term extreme wildfire event appearing four times in total; the beneficial role of fire is absent, and nature-based solutions appear only marginally in the forestry and disaster-risk clusters despite their prominence in climate and biodiversity policy \citep{plana_framing_2024}. A comparison of European and United States arrangements for critical infrastructure reaches a parallel conclusion, identifying insufficient integration of scientific evidence into policy and recommending standardised cross-border risk assessment \citep{kalapodis_towards_2025}. The Alpine coordination problem is therefore not only a matter of eight countries failing to align with one another; the supranational framework they would align to barely names the practice.

Harmonisation has been attempted, but it has been project-driven. The first harmonised Alpine fire database and regime characterisation came from a transnational project \citep{valese_wildfires_2011}, whose data layers were subsequently reused in national danger assessment \citep{muller_towards_2020}, and whose synthesis was consolidated into an Alpine-scale white paper setting out prevention, suppression, post-fire and knowledge-transfer priorities with costed estimates \citep{muller_forest_2020}. What does not yet exist at Alpine scale is standing infrastructure: the multi-institutional monitoring networks combining remote sensing, in situ instrumentation and geophysical methods that compound hazard research requires \citep{baselt_compound_2026}, and the systematic post-fire documentation procedures that would make indicator-based integrated fire management operable \citep{melzner_wildfire-induced_2025}. That fire and its consequences cross national borders has been demonstrated directly for a central European transboundary interface \citep{avouris_exposure_2025}.

Standing coordination across the jurisdictions of a mountain range does exist elsewhere, and the one documented case shows both what it can carry and where it stops. The protected areas of the Australian Alps fall under the legislation of three states and territories, and since 1986 a memorandum of understanding between four governments has supported a cooperative management programme coordinated by the Australian Alps Liaison Committee. The programme treats the range as one biogeographic region for fire management, water and catchments, invasive species and Aboriginal engagement, and it has survived political and funding cycles and major fire events \citep{lockwood_biodiversity_2014}. Coordination did not insulate fire policy from politics, however. After the 2009 Victorian fires a hazard-reduction burning target of 5\% of public land per year became state policy, although such burning is not undertaken in the treeless alpine landscape. In 2010 cattle were returned to the Victorian Alps for a trial of their contribution to reducing fire risk, despite published evidence suggesting that alpine grazing has no significant effect on it, until the federal government removed them on the strength of the range's national heritage listing \citep{lockwood_biodiversity_2014}. The case is single and imported, but it suggests that a standing body can align agencies without settling the land-use disputes that fuel management raises.

\subsection{The economics of Alpine wildfire}
\label{sec:economics}

The economics of Alpine wildfire is ordered counter-intuitively, since what the forest protects exceeds what it produces and what prevention costs is small beside both.

Protective function dominates the value at risk. Replacement-cost valuation of a ten-hectare Alpine Scots pine protection stand returns €30{,}440 per hectare, or €950 per hectare per year \citep{accastello_asforesee_2019}. Avoided-risk valuation of a Valais protection forest places annual rockfall risk at 19{,}100\,CHF with the forest and 95{,}500\,CHF without it. Under moderate and high emissions scenarios those figures rise to 59{,}250 and 70{,}100\,CHF per year within 150 years, as composition shifts and small blocks reaching the railway move from decadal to triennial recurrence \citep{moos_climate_2021}. Neither figure is a timber value, and in terrain where extraction is technically demanding and economically marginal \citep{dinca_timber_2025} the timber value is frequently the smaller number.

On the expenditure side, response dominates. Direct firefighting and post-fire management across the Alpine arc are estimated at around €75 million annually, against a proposed integrated fire management framework costed at roughly €10 million a year, of which €5--7 million would go to prevention, €1.5--2.5 million to suppression capability, €1--2 million to post-fire management and €0.5 million to knowledge transfer \citep{muller_forest_2020}. The ratio, whatever the precision of the estimates, frames prevention as inexpensive relative to the response it would displace. Econometric analysis of Italian regions reaches the same conclusion from a different direction, concluding that upstream investment in prevention is more efficient than suppression under projected climate intensification \citep{michetti_forest_2019}.

That analysis also supplies the clearest quantitative statement of how socio-economic and climatic drivers combine in a European mountain setting. Across nineteen Italian regions, a ten-millimetre increase in monthly precipitation reduced fire frequency by 8.9\% in the same month with effects persisting longer in the north than the south, while a one-degree temperature increase raised frequency by up to 15\% in the south and between 0.9\% and 5.7\% elsewhere. Socio-economic terms were of comparable magnitude: a one-percentage-point rise in tertiary education reduced fire events by 11\% in the north, and in the south a single additional extortion per hundred thousand inhabitants raised fire frequency by 4.6\% and burnt area by 7\%. Projected to 2016--2035, the northern cluster gains 11--12\% in annual fire events and 23--27\% in annual burnt area \citep{michetti_forest_2019}. Education and institutional quality are fire management variables, and they are not typically treated as such.

Delivery is a further problem, and one that fragmented Alpine forest ownership shares with the case where it has been measured. In Catalonia, wildfire prevention subsidies funded by a vehicle carbon tax and offering €5{,}000 to €200{,}000 per proposal reached only 3.87\% of the 151{,}196\,ha designated as strategic, and the average minimum distance between subsidised parcels was 287\,m, so the spatial continuity the scheme was designed to create was not achieved; the barriers identified were non-adjacent parcels, uneven social capital between owner associations, and reimbursement rules requiring owners to advance costs \citep{gorrizmifsud_current_2026}. Fuel treatment that must be contiguous to work is being funded through instruments that allocate to whoever applies.

The economics literature also identifies structural incentive problems that the Alpine studies have not yet examined. A private parcel owner deciding on mitigation weighs private costs and benefits while ignoring spatial spillover externalities (fire transmitted across parcel boundaries) and conventional externalities such as smoke; liability rules and insurance markets then shape mitigation incentives in ways that vary by jurisdiction \citep{bayham_economics_2022}. In an Alpine context of small private parcels and eight national legal systems, these are first-order determinants of whether fuel treatment happens at all, and the Catalan subsidy experience is a direct illustration of the spillover problem going unpriced \citep{gorrizmifsud_current_2026}.

Pastoral systems sit on both sides of the Alpine ledger and neither side has been priced. Grazing is the fuel-treatment instrument with the broadest policy leverage (Section~\ref{sec:forest}) and is inexpensive relative to mechanical clearance because the work is performed by animals kept for other reasons \citep{xu_how_2026}, so its marginal cost to a fire budget is largely the cost of sustaining a pastoral economy that already exists and is contracting. On the other side, alpine pasture, its access infrastructure and the stock on it are exposed assets, and transhumant animals are at their highest elevation during the summer fire season. We found no Alpine study valuing either side: no estimate of the fuel-management service that extensive grazing supplies without being paid for it, and no record of livestock or pasture loss to fire against which to set it. Since the agricultural-landscape component carried the entire societal-cost saving in the one European landscape optimisation available \citep{lecinadiaz_incorporating_2023}, this is a quantitative gap with direct consequences for how agri-environmental payment is justified.

What the economics cannot yet resolve, however, is distribution. Alpine valley communities are small, ageing and economically specialised; they bear the costs of both the hazard and its prevention, while a large share of the protective benefit accrues to through-traffic, downstream water users and regions outside the mountains. The services at stake propagate across scales, since the vegetation changes through which fire acts on them operate locally but register regionally \citep{vukomanovic_systematic_2019}, and risk assessment frameworks that make adaptive capacity an explicit component of vulnerability come closest to representing this asymmetry \citep{riviere_mapping_2023, lecinadiaz_assessing_2021}. Tourism sits on both sides of the ledger and its seasonal redistribution moves exposure into the fire season, as warming reduces high-altitude snow-cover duration while drawing summer visitors escaping lowland heat \citep{pede_mountain_2022, abrham_wildfires_2024}. We found no Alpine study that quantifies who bears the costs and who receives the benefits.

\subsection{What differs in mountains}

For Alpine managers, the importance of a fire depends partly on the protective function it affects, even when the burnt area is small \citep{muller_forest_2020}. Response planning therefore needs information about assets below the fire as well as expected spread. The same terrain can constrain crew access, water delivery and safe working positions, making coordination with mountain rescue and aerial resources important.

Access also conditions prevention. Mediterranean evidence shows that fuelbreak containment is much more successful when ground crews can work on the line, while a multi-region analysis links difficult terrain to lower initial-attack success \citep{ortega_modeling_2024, cardil_assessing_2025}. These results justify evaluating access in Alpine treatment design. Their numerical effects cannot be adopted without testing under local tactics and resources.

Beyond individual incidents, the Alpine Convention perimeter spans eight countries with different reporting and management systems. Harmonisation has relied substantially on collaborative projects \citep{valese_wildfires_2011}. Fragmented ownership can further complicate coordinated treatment, a problem also documented in Catalonia \citep{gorrizmifsud_current_2026}. Standing cooperation in the Australian Alps offers one institutional comparison, rather than a ready-made model for the European range \citep{lockwood_biodiversity_2014}.

Coordination also requires agreement on who pays. Forest protection benefits settlements and infrastructure, while treatment costs fall on particular owners and authorities \citep{accastello_asforesee_2019, moos_climate_2021}. Italian evidence associates active land governance with lower fire impacts, but does not establish the marginal return on Alpine prevention spending \citep{spadoni_active_2023}. Cost-effectiveness and the distribution of benefits therefore remain questions for evaluation, rather than conclusions that follow from the current prevention and suppression budgets.

% =======================================================
\section{Discussion}
\label{sec:discussion}

The evidence supports a mountain-specific approach to wildfire assessment, but not a single mountain fire regime. Relief affects local weather, spread, access and the routes by which damage reaches exposed assets. Ignition patterns, fuel histories and institutional arrangements vary between ranges and within the Alps. Table~\ref{tab:contrast} summarises the implications without assuming a uniform Mediterranean or lowland baseline.

\begin{table}[H]
\centering
\footnotesize
\setlength{\tabcolsep}{4pt}
\renewcommand{\arraystretch}{1.1}
\caption{Implications of the reviewed mountain evidence for assessment and management. These are conditional findings, not universal contrasts between mountain and lowland fire.}
\label{tab:contrast}
\begin{tabularx}{\textwidth}{@{}>{\raggedright\arraybackslash}p{2.1cm} >{\raggedright\arraybackslash}X >{\raggedright\arraybackslash}X >{\raggedright\arraybackslash}p{3.0cm}@{}}
\toprule
\textbf{Domain} & \textbf{Relevant finding} & \textbf{Assessment implication} & \textbf{References} \\
\midrule
Ignition & Human and lightning ignitions differ in season and terrain & Test separate seasonal and ignition-source relationships & \citep{vacchiano_modeling_2018, conedera_lightninginduced_2006} \\
Spread & Wind--slope interactions and model assumptions affect predicted behaviour & Validate spread rate and intensity as well as perimeter & \citep{malanchini_modeling_2025, hayajneh_wind_2026} \\
Observation & Many small fires are absent from coarse products; phenology affects severity retrieval & Combine event registers, suitable imagery and field observations & \citep{bakke_datadriven_2023, morresi_mapping_2022} \\
Climate & Regional trends coexist with strong local weather gradients & Evaluate conditions at the elevations and slopes exposed to fire & \citep{wastl_recent_2012, schunk_forest_2013} \\
Water and slopes & Fire changes runoff, sediment supply and forest protection on different timescales & Use process-specific thresholds and recovery trajectories & \citep{rengers_evaluating_2024, gehring_shallow_2019} \\
Forest management & Treatment changes fuels and potentially protective structure & Evaluate species response, retained structure and downslope assets together & \citep{maringer_temporal_2016, gamba_prescribed_2025} \\
Risk & Direct and cascading losses can extend beyond the burnt area & Link hazard to exposure, vulnerability and recovery & \citep{laporta_postwildfire_2025, papathomakohle_towards_2025} \\
Governance & Access, ownership and responsibilities constrain implementation & Coordinate prevention, suppression and post-fire management across authorities & \citep{kirschner_governance_2024, gorrizmifsud_current_2026} \\
\bottomrule
\end{tabularx}
\end{table}

For the Alps, the practical concern is the interaction between changing fire conditions and forests that protect settlements and infrastructure. Projected hazard increases coincide with changes in forest extent and continuity, but neither trend alone determines future losses \citep{gerberding_mapping_2025, garbarino_contrasting_2020}. Location, fire severity, retained protective structure and subsequent weather all matter. We therefore hesitate to infer regional damage from burnt area alone.

This emphasis on consequences does not establish that post-fire geomorphic hazards dominate Alpine losses. Direct damage, smoke, ecological change and disruption of transport or tourism remain relevant. The arc-wide regime assessment placed current fire damage below that from storms and bark beetles \citep{conedera_characterizing_2018}. The reviewed evidence does not overturn that comparison. It identifies pathways through which a relatively small fire could nevertheless become locally consequential.

The modelling implication is to match each tool to its decision. An occurrence map can support prevention targeting, but cannot determine spread intensity, structural damage or a debris-flow runout. Conversely, a detailed event model does not establish the regional frequency of the event it simulates. Model development should therefore combine clearly defined outputs, representative observations and validation at the scale of use. Comparative studies provide mechanisms and candidate methods; Alpine evaluation determines which numerical relationships can be adopted.

Improved observation is central to that evaluation. Small-fire records are needed to assess occurrence, while post-fire monitoring should follow soil properties, sediment mobilisation and forest protection separately \citep{melzner_wildfire-induced_2025, baselt_compound_2026}. The objective is to establish how each process changes after fire and when its influence becomes negligible for a particular asset. Linking these observations to exposure and damage would also allow vulnerability assumptions to be tested, rather than carried forward solely from expert judgement.

Management trials can address a complementary set of questions. Fuel treatment in protection forests should be evaluated against both fire behaviour and retained protective capacity. The single Alpine prescribed-burning trial supports further work across species, sites and longer observation periods \citep{gamba_prescribed_2025}. Grazing and mechanical treatment also need local evaluation, including regeneration and erosion effects \citep{rouetleduc_effects_2021, xu_how_2026}. For linear treatments, access and crew safety are essential design criteria because containment depends strongly on suppression at the line \citep{ortega_modeling_2024}. Deadwood retention and salvage require the same attention to site conditions and the assets below.

The value of those interventions cannot be inferred from current expenditure alone. Prevention and response budgets describe what is spent, not the loss avoided by an additional unit of either activity. Existing valuation methods could support comparisons that include forest protection, treatment costs and the distribution of benefits \citep{accastello_asforesee_2019, moos_climate_2021}. Such analysis would help identify when landowners, municipalities or infrastructure operators have different incentives, and where shared funding is justified.

These decisions extend beyond fire services. The Aran valley programme and cooperation in the Australian Alps show how shared objectives can support management across institutional boundaries \citep{oliveres_path_2025, lockwood_biodiversity_2014}. They do not establish that the same arrangements will work in the European Alps. We propose clearer links between prevention, suppression and post-fire responsibilities, with a fire triggering joint assessment by forest, fire and natural-hazard authorities. This is consistent with place-based adaptation and coexistence approaches \citep{moritz_learning_2014}, while recognising the Alps' existing suppression institutions and experience.

Several consequences remain insufficiently measured. These include the combined effects of fire on Alpine snow and glacier-fed water supply, faunal mortality and longer-term ecological change. Evidence elsewhere shows that some effects compete: deposited material can increase ice melt while atmospheric smoke reduces incoming radiation \citep{aubrywake_fire_2022}. These interactions require measurement before they are incorporated as directional feedbacks in Alpine scenarios.

\label{sec:limitations}
The review also has methodological limits. Keyword searches and citation chaining do not constitute an exhaustive database search, and the starting literature influences what is found. The available search record specifies term combinations but does not preserve a complete dated export of Google Scholar results. Selection informed by the author team's disciplines may also favour some topics over others. Targeted supplementation improves coverage but cannot demonstrate that no relevant literature has been missed.

Language, indexing and access impose further constraints. Open citation infrastructure incompletely represents national reports and other grey literature, while full-text availability influences which candidates can be appraised. This matters in a range where important operational knowledge is published in several national languages. The disciplinary review loop can identify omissions, but remains open in this draft and should not be presented as a completed validation of coverage.

Heterogeneous event definitions, thresholds and study designs prevent quantitative pooling. We use reported values to characterise particular settings, not as interchangeable estimates of an Alpine parameter. The Hallstatt rockfall case, the Bussoleno debris flow and the long-term subalpine regeneration study each concern a single event or site \citep{melzner_brief_2019, laporta_postwildfire_2025, degiuli_arrested_2024}. They establish that these responses are possible, while leaving their regional prevalence uncertain.

Comparative evidence helps assess mechanisms but cannot replace local data. Similar outcomes across ranges may arise under different fuels, weather and management, and agreement alone does not prove a mountain-specific cause. Cross-range synthesis is useful when those conditions and limits remain explicit \citep{payne_nature_2020}. We retain that distinction between an observed process, a plausible Alpine application and an Alpine parameter estimate.

% =======================================================
\section{Conclusions}
\label{sec:conclusions}

Wildfire assessment in the European Alps needs to account for the interaction of terrain, land use and forest protection. Most recorded fires are small, but their consequences depend on what burns and what lies below it. The evidence supports considering direct damage together with erosion, debris flow, rockfall and changes in protective forest structure.

These pathways recover at different rates. Soil responses measured over several years cannot be equated with losses of root reinforcement or stem protection lasting decades. Likewise, persistent failure of tree regeneration at one site does not establish permanent elevation of every post-fire hazard. Risk assessment needs recovery trajectories for the process and asset concerned.

The available models provide useful components of that assessment. Their transfer requires attention to local weather, fuels, ignition regimes and the observations used for validation. Better small-fire records and sustained post-fire monitoring would improve those tests, while exposure and vulnerability information would connect hazard estimates to potential loss.

Management must make the same connection. Fuel reduction, forest restoration and land-use policy should be assessed alongside the protection that forests provide to settlements and infrastructure. The immediate priority is to evaluate these trade-offs in Alpine settings and coordinate responsibility through the post-fire period. Evidence from other ranges offers methods and practical precedents, while local observation and disciplinary review remain necessary to establish their suitability.


\section*{Acknowledgements}
\addcontentsline{toc}{section}{Acknowledgements}
% Text here.

\section*{Funding}
\addcontentsline{toc}{section}{Funding}
% Text here.

\section*{Competing Interests}
\addcontentsline{toc}{section}{Competing Interests}
% Text here.

\section*{Data and Code Availability}
\addcontentsline{toc}{section}{Data and Code Availability}
% The screening register underlying this review (all candidate records, tier,
% inclusion decision and reason) is provided as supplementary material.


% =======================================================
\appendix
% =======================================================

\section{Search Protocol and Screening Register}
\label{app:protocol}

\subsection{Counts at each step}

The six steps of Section~\ref{sec:scope} and Figure~\ref{fig:flow} are stated below with the counts each produced. The combination of starting documents, keyword search and citation chaining suits a literature that is administratively fragmented, multilingual and heavily dependent on national technical reporting, which is the condition under which database querying performs least well \citep{mantero_influence_2020}.

\textbf{Step 1, starting documents.} The EUSALP white paper \emph{Forest fires in the Alps: state of knowledge, future challenges and options for an integrated fire management} \citep{muller_forest_2020} covers the whole Alpine Convention perimeter. It combines harmonised fire statistics for the arc, a survey of authorities, scientists and action forces across the Alpine countries, and costed recommendations for prevention, suppression, post-fire management and knowledge transfer, and its reference list holds 120 publications and web sources together with six national fire databases. The Eurac Research co-authored study \citep{gerberding_mapping_2025} maps the probability of forest fire hazard across the arc under climate change scenarios.

\textbf{Step 2, keyword search and seed library.} The Google Scholar searches used the term combinations wildfire, mountain and Alps, and then wildfire and mountains. Candidates were retained where they bore on a discipline represented in the in-house expertise of Eurac Research, namely natural hazards and risk, climate, drought and water, forest science, and economics and governance, which are the disciplines of the co-authors who review the sections in Step~6. The seed library holds 49 distinct publications, including the Eurac Research study; 31 appeared in or after 2020, three also stand in the white paper's reference list and six cite the white paper. The library as supplied contained 100 bibliographic entries, because 44 records had been exported more than once, three of them four times; this repetition is an artefact of how the library was compiled and is reported here for transparency rather than as a screening result. Of the 49 seed publications, 47 were available as full-text documents.

\textbf{Step 3, citation chaining.} The reference sections of the white paper and of 45 seed documents were extracted directly from the source files, yielding 3{,}575 reference records, 120 of them from the white paper. Matching on digital object identifier and on normalised title reduced these to 3{,}375 distinct works. Of the white paper's 120 references, 62 were recovered independently through the seed documents, and 58 entered the register through the white paper alone. For forward chaining, the 14 seeds most central to the Alpine question were submitted to the OpenCitations Index and Meta, returning 508 citing works, of which 319 resolved to usable metadata within this iteration. The remaining 189 are recorded in the register for a subsequent pass.

\textbf{Step 4, targeted additions.} This step is the reason the corpus is weighted towards open-access and recent publication, and it is the step least reproducible from the seed library alone.

\textbf{Step 5, screening.} Cross-source de-duplication left 3{,}228 records for screening. Every record was assigned to one of four tiers, with a recorded reason: \emph{include} (fire content and the European Alps or an Alpine Convention country); \emph{comparative} (fire content and another mountain system, admitted on the strength of the process rather than only where the Alps are silent, following the two-tier rule of Section~\ref{sec:scope}); \emph{context} (mountain process without fire, or fire without a mountain setting); \emph{exclude} (generic fire science with no mountain anchor). Screening excluded 2{,}332 records and retained 896, comprising 290 include, 194 comparative and 412 context. Full text was obtained and read for 136 of the retained records, prioritised by their mapping to the review's thematic sections, and 150 sources are cited in the text.

\textbf{Step 6, discipline review.} Climatologists review the climate and meteorology evidence of Section~\ref{sec:climate}, drought and hydrology researchers the drought and water evidence of Section~\ref{sec:water}, and forest scientists the forest evidence of Section~\ref{sec:forest}. Natural-hazard and risk specialists review Sections~\ref{sec:hazard} and~\ref{sec:risk}, and economists and governance researchers review Section~\ref{sec:governance}.

\subsection{What was recorded for each source}

Each screened record carries a tier, a decision, a reason and a mapping to one or more sections of this review. For every source cited in the text, a structured evidence note recording study region, method, quantitative results and stated limitations was captured and archived alongside the screening register. Both are provided as supplementary material. Where a claim in the text rests on fewer than three independent cases, the number is stated in the text, and mechanisms imported from outside the Alps are identified as imported at the point of use.

Synthesis is thematic rather than quantitative, for the reasons of definitional and administrative heterogeneity set out in Section~\ref{sec:scope}.

\subsection{What the method determines about the corpus}

The method shapes what the corpus contains, and its effects should be read as constraints on coverage rather than on the reliability of individual sources. Citation chaining reproduces the coverage of its starting point, so topics well represented in the white paper and among the seeds are recovered densely and topics absent from both are recovered weakly. Forward chaining operates on open citation infrastructure, which indexes journal literature comprehensively but national technical reporting and non-English publication sparsely, a material consideration for a range whose fire statistics were harmonised only recently \citep{valese_wildfires_2011, muller_forest_2020}. The expansion phase was confined to material whose full text could be retrieved, which weights the added corpus towards open-access and recent publication.

A subsequent iteration should add database searching with recorded strings and dates, resolve the 189 forward records held over, and target the national fire and forest reporting of the Alpine Convention countries directly, where much of the operational and grey-literature knowledge base resides.

\section{Glossary of Terms}

\begin{description}
    \item[Fire danger] The meteorological and fuel-condition potential for ignition and spread, expressed by an index; a statement about conditions, not about realised fire.
    \item[Fire regime] The realised distribution of fire frequency, size, severity, seasonality and cause over a region and period.
    \item[Burn severity] The degree of fire-induced change to vegetation and soil, commonly estimated remotely; see Section~\ref{sec:water} for the limits of this proxy.
    \item[Fuel connectivity] The extent and spatial aggregation of contiguous burnable vegetation, distinct from fuel load.
    \item[Protection forest] Forest whose primary managed function is mitigation of gravitational natural hazards affecting assets below.
    \item[Geohydrological process] Collective term used in the Alpine literature for erosion, rockfall, slides and debris flows driven by the interaction of water and slope material.
    \item[Cascade] A sequence in which one hazard alters the susceptibility of a system to a subsequent, distinct hazard.
    \item[Firebreak] A linear strip from which fuel is removed entirely, intended to stop fire spread.
    \item[Fuel break (\emph{coupure de combustible})] A linear strip in which fuel load and vertical continuity are reduced but the stand is retained; functions primarily as a platform for ground suppression rather than as a barrier.
    \item[Exposure] The elements (people, buildings, infrastructure, habitats, protective forest, livestock and pasture) located where a hazard and its consequences can reach them.
    \item[Vulnerability] The susceptibility of an exposed element to loss given a hazard of stated intensity; distinct from exposure and from hazard, and measured by damage-calibrated or judgement-weighted indicators.
    \item[Ecological vulnerability] Potential loss of ecological value integrated over the post-fire recovery period, combining the value of what burns, the resistance of the vegetation to fire and the time taken to return to pre-fire condition.
\end{description}


% =======================================================
% Bibliography
% =======================================================
\bibliographystyle{apalike}
\bibliography{FireScape_master_v6}

\end{document}
